\textbf{\textsc{MASTER\textquotesingle S THESIS - ENGLISH REVIEW DRAFT}}

Auditable Governance of Long-Term Memory in Language Agents

The ASK Controller (Keep-Supersede-Abstain), the MGP-6 v2 Diagnostic,
and Evidence-Bounded Evaluation

Review state: 15 July 2026\\
Evidence basis: ledger v1.1 + construct-validity audit\\
Editable research draft

\begin{quote}
\textbf{The official University of Tehran template and verified
administrative fields were not supplied. This document is a review copy,
not an institution-ready submission.}
\end{quote}

Administrative fields intentionally remain external

\section{Front Matter}\label{front-matter}

\textbf{Master\textquotesingle s Thesis --- English Review Draft}
Administrative fields (candidate, supervisor, advisor, faculty, dates,
signatures) are intentionally left as placeholders and are consolidated
in frontmatter/ADMINISTRATIVE\_PLACEHOLDERS.md. They must never be
filled in by tooling.

Persian-language front matter is available in
frontmatter/abstract\_fa.md, frontmatter/extended\_summary\_fa.md, and
frontmatter/terminology\_fa.md.

\textbf{Evidence-status notice.} Every number in this thesis is
transcribed from the frozen evidence ledger
artifacts/results\_ledger.json v1.1 (sha256
fea886f1ae746e921e785536e6ce47bea89ec11b99f8f2e75936dc71c6c5185e,
created 2026-07-13) or the later construct-validity audit
artifacts/analysis/mgp6\_validity\_audit.json. The 2026-07-15 shortcut
failure supersedes the ledger\textquotesingle s older final\_candidate
interpretation for MGP-6 v2: those policy rows are now diagnostic. No
failure\_replay\_only, incomplete, invalid, externally\_pending, or
mapping-dependent row is used as a primary capability result. Conversion
to the official University of Tehran thesis template and completion of
its administrative fields remain external requirements (see
thesis/README.md and thesis/thesis\_checklist.md).

\subsection{Declaration}\label{declaration}

This manuscript distinguishes verified results, supported-with-caveat
results, diagnostic observations, and hypothesized designs, following
the strength tags of reports/CLAIMS\_EVIDENCE\_MATRIX.md. Invalid and
superseded runs are named and excluded
(docs/audit/INVALID\_AND\_SUPERSEDED\_RUNS.md). Raw and adjudicated
numbers are always reported together. No LongMemEval number is presented
as final.

\subsection{Abstract}\label{abstract}

Long-term memory makes language agents more useful and, at the same
time, creates a governance problem: the agent must decide what to keep,
what to supersede when new evidence contradicts stored state, and when
to abstain because the evidence is insufficient to justify a write.
These write-time decisions are not the same construct as downstream
answer accuracy, and evaluations that collapse them conceal both
capability and risk.

This thesis formulates memory governance as an explicit, auditable
three-way decision --- KEEP, SUPERSEDE, ABSTAIN --- over typed,
provenance-tracked records, and separates evaluation of the write
action, the resulting state and lineage, and the later answer. It
reports two instruments. The first, MGP-6 v2, is an author-designed
180-item pilot whose most important result is a \emph{failed}
construct-validity gate: a conflict-type-only lookup fitted on train and
development reaches 1.000 held-out accuracy without reading the dialogue
or memory state (type-only: 0.9667). That failure motivated the second
instrument.

MGP-6 v4.3 is a 900-item human-authored controlled synthetic diagnostic
in 100 semantic components across 20 domains, split by component into
exactly action-balanced train (180), development (180), and sealed
hidden (540) partitions, with duplicate, shortcut, allowlist, and
deterministic-build gates passed at release and dual human annotation
covering every item. One pre-registered configuration --- a composed
controller pairing a public-evidence rule policy with a live model
policy under a conservative abstention guard --- was run exactly once on
the sealed split. It attains \textbf{0.8481 accuracy} and \textbf{0.8479
macro-F1} (component-cluster 95\% intervals {[}0.8037, 0.8889{]} and
{[}0.8021, 0.8894{]}; Cohen\textquotesingle s κ = 0.7722 against
balanced gold). The component-cluster interval is roughly 40\% wider
than an item-independence interval, which is the correct width and the
one reported.

The thesis argues that the error \emph{structure} carries more
governance information than the aggregate, and reports it accordingly.
Of 38 mishandled ambiguous items, 30 become KEEP --- leaving the stored
state byte-identical and losing only the audit record --- while just 8
replace an active record on insufficient evidence (0.0444, CI {[}0.0227,
0.0852{]}). A single "false commit" rate of 0.2111 therefore overstates
the state-corrupting harm roughly fivefold. Missed updates run 0.0833
and destructive overwrites 0.0611. Component decomposition shows neither
the rule policy nor the model policy alone exceeds 0.750 on development
data, so the result is a genuine composition rather than a rule engine
with a model attached. Each write decision costs a 1.3 s median and 393
prompt tokens, below every full-context baseline it outperforms;
monetary cost is recorded as unavailable rather than estimated, because
provider telemetry emitted no price field.

On the development split, every structural baseline, every component
ablation, the zero-shot policy, the standalone rule policy, and three
deployed external memory systems fall between 0.2944 and 0.3667 --- at
or below the 0.3333 chance rate of the balanced decision --- failing in
two opposite ways: commit-happy systems overwrite 77--98\% of records
they should keep, while abstain-happy systems decline 78--98\% of
decisions they should make. The composed controller reaches 0.8444 and
its successor configuration 0.9611, with every
controller-versus-baseline contrast rejecting at Holm-corrected p
\textless{} 1.6e−19.

As downstream evidence, a homogeneous full-slice LongMemEval-M run (all
180 items, one answer model, one judge, zero judge fallbacks) places the
governed controller at 76.1\%, above every structural baseline (naive
overwrite 63.3\%; paired p = 2.9e−04), with the gain concentrated in
knowledge-update and preference competencies rather than spread
uniformly --- the pattern a write-policy effect should produce.

What the evidence cannot support is stated as flatly as what it can.
Only one hidden configuration was authorised, so every baseline and
ablation contrast is development evidence on the split used for
selection; the collections are synthetic, English-only, single-slot, and
single-step; human annotation is owner-attested rather than
identity-authenticated; and the abstention mechanism is not a calibrated
selective predictor --- its selective accuracy (0.8449) falls below its
overall accuracy (0.8481). The thesis contributes an auditable
evaluation decomposition, a released benchmark with build-time construct
gates, per-item predictions with cluster-aware and cost-conditional
analysis, and a prospective independently generated multi-human
confirmation protocol. It does not establish population safety,
calibrated confidence, or long-lived deployment performance.

\textbf{Keywords:} long-term memory, language agents, memory governance,
supersession, abstention, provenance, lineage, auditability, diagnostic
evaluation, construct validity, selective prediction, error
decomposition, MGP-6 v4.3, ASK Controller, LongMemEval.

\subsection{Acknowledgements}\label{acknowledgements}

Institutional and personal acknowledgements are administrative front
matter and will be added from
frontmatter/ADMINISTRATIVE\_PLACEHOLDERS.md before submission. The
technical package acknowledges the LongMemEval authors (Wu et al.,
2025), the maintainers of the evaluated external memory systems, and the
annotation panel configuration documented in the ledger.

\subsection{Reader\textquotesingle s guide}\label{readers-guide}

Chapter 1 frames the research problem and questions. Chapter 2 situates
the work in agent memory systems, memory lifecycle research, benchmarks,
belief revision, and selective prediction, including the 2026 near
neighbors. Chapter 3 formalizes the problem and threat model. Chapter 4
specifies the ASK Controller as implemented. Chapter 5 documents MGP-6
v2 and its annotation evidence. Chapter 6 defines the experimental
methodology and the claim-eligibility policy. Chapter 7 reports results.
Chapter 8 gives the error analysis and discussion. Chapter 9 addresses
validity, ethics, and limitations. Chapter 10 concludes. Appendices A--F
provide schemas, prompt pointers, metric definitions, the
annotation-manual pointer, reproducibility commands, and the artifact
manifest.

\subsection{Abbreviations}\label{abbreviations}

{\def\LTcaptype{none} % do not increment counter
\begin{longtable}[]{@{}
  >{\raggedright\arraybackslash}p{(\linewidth - 2\tabcolsep) * \real{0.2610}}
  >{\raggedright\arraybackslash}p{(\linewidth - 2\tabcolsep) * \real{0.7390}}@{}}
\toprule\noalign{}
\begin{minipage}[b]{\linewidth}\raggedright
\textbf{Term}
\end{minipage} & \begin{minipage}[b]{\linewidth}\centering
\textbf{Meaning}
\end{minipage} \\
\midrule\noalign{}
\endhead
\bottomrule\noalign{}
\endlastfoot
ASK & Keep--Supersede--Abstain: the three-action memory-governance
controller of this thesis \\
MGP-6 v2 & Author-designed six-type memory-governance diagnostic slice
(180 scenarios) \\
LME-M & Stratified 180-item LongMemEval-M evaluation slice (6
competencies × 30, seed 0) \\
SSP & Single-session preference (LongMemEval competency) \\
RD & Paired risk difference (accuracy difference on the same items) \\
CI & Confidence interval \\
IAA & Inter-annotator agreement (not claimed here; see §5.5) \\
n.s. & Not statistically significant \\
JSONL & JSON Lines format \\
\end{longtable}
}

\section{1. Introduction and Research
Questions}\label{introduction-and-research-questions}

\subsection{1.1 Motivation}\label{motivation}

A language model with no persistent memory cannot reliably carry
preferences, commitments, or corrections across sessions. A language
model with persistent memory is risky for the opposite reason: stale
facts can keep influencing behavior after the user has corrected them,
contradictory records can accumulate silently, and a system can act on
information it should have refused to write in the first place. The
engineering objective is therefore not "remember as much as possible"
but to govern the memory write as a decision under explicit policy.

This observation changes where control must live. Prompting a model to
"be careful" is insufficient because memory writes have side effects
that outlive the prompt. A new candidate fact must be classified against
the existing store --- is it independent, a duplicate, or a
contradiction of a prior record? --- and the resulting action must be
explicit and auditable: keep the new record alongside the old, supersede
the old record while preserving its lineage, or abstain because the
evidence is too ambiguous to commit either way. When a system cannot
express the third option, every conflict becomes a silent overwrite or a
silent duplicate.

Existing evaluations mostly treat memory as an answer-retrieval problem:
given a long interaction history, can the model answer a later question
(Wu et al., 2025; Maharana et al., 2024)? That view is valuable but
incomplete: end-to-end question accuracy cannot attribute a failure to
the write decision that caused it, and it cannot see governance failures
that happen not to affect the sampled questions. A governance evaluation
must score the decision boundary itself.

\subsection{1.2 Research problem}\label{research-problem}

How can a language-agent memory system make and audit policy-sensitive
write decisions --- keep, supersede, abstain --- and how can those
decisions be evaluated without conflating governance correctness with
retrieval or answer quality?

The problem decomposes into five concerns: \textbf{representation}
(typed memory records with provenance and lineage), \textbf{decision}
(an interpretable, abstention-aware action policy), \textbf{execution}
(supersession must actually mark the prior record inactive and preserve
its lineage), \textbf{measurement} (decision-level metrics with correct
denominators, plus state metrics such as lineage integrity and stale
retrieval, distinct from answer metrics), and \textbf{honest reporting}
(every number carries a claim-eligibility label; contaminated or patched
runs are labeled as such).

\subsection{1.3 Research questions}\label{research-questions}

\textbf{RQ1 --- Architecture.} What controller decomposition makes
memory-write decisions explicit, interpretable, and auditable, and
prevents a language model from mutating memory outside a governed action
vocabulary?

\textbf{RQ2 --- Diagnostic.} Can a compact, author-designed diagnostic
slice expose memory-governance failures --- invalid commits, premise
acceptance, lineage loss --- that are invisible to downstream answer
accuracy?

\textbf{RQ3 --- Measurement.} Which metrics separate decision
correctness, memory-state correctness, and downstream answer behavior,
and which shipped metrics fail (and must be disclosed as failing)?

\textbf{RQ4 --- Reliability.} How far do the reported conclusions
survive paired significance testing with multiplicity control,
annotation audit, and adversarial review of run provenance?

\textbf{RQ5 --- Reproducibility.} What artifact and claim controls keep
a thesis, its papers, and its code mutually consistent --- including
when the honest answer is "this number is not final"?

\subsection{1.4 Hypotheses and their outcomes
(preview)}\label{hypotheses-and-their-outcomes-preview}

The original directional comparisons are not confirmatory because the
pilot fails its shortcut gate. Chapter 7 therefore treats them as
diagnostics and reports these outcomes:

\begin{itemize}
\item
  \textbf{H1 --- Governance framing helps an LLM.} Not identified by v2:
  structured and unstructured prompts differ on a label-leaking
  collection, so the observed gap is descriptive.
\item
  \textbf{H2 --- An interpretable reference policy can match an LLM
  policy.} No superiority or non-inferiority conclusion: test accuracy
  is 0.8889 versus 0.9000 and the family-clustered paired interval
  includes zero.
\item
  \textbf{H3 --- Abstention is load-bearing for action safety.}
  Mechanistically supported inside the generator, but not validated
  externally; zero observed action false commits has a non-zero
  finite-sample upper bound and accompanies missed updates.
\item
  \textbf{H4 --- Decision and stored-state metrics can diverge.}
  Supported as a measurement demonstration: the no-lineage ablation
  preserves action accuracy while degrading lineage integrity.
\item
  \textbf{H5 --- Controller end-task superiority on LongMemEval.} Not
  supported; existing downstream runs are contaminated, incomplete, or
  selection-biased.
\end{itemize}

\subsection{1.5 Contributions}\label{contributions}

\begin{enumerate}
\def\labelenumi{\arabic{enumi}.}
\item
  \textbf{Formulation and evaluation decomposition (RQ1, RQ3).} Memory
  writes are represented as KEEP/SUPERSEDE/ABSTAIN decisions over typed,
  provenance-tracked, lineage-preserving records and evaluated
  separately at action, state, and answer layers.
\item
  \textbf{A benchmark with build-time construct gates (RQ2, RQ4).} MGP-6
  v4.3 --- 900 human-authored items in 100 semantic components across 20
  domains, component-split, exactly action-balanced, dual-annotated,
  with shortcut/duplicate/allowlist/deterministic-build gates enforced
  at release rather than audited afterwards, and a sealed 540-item
  hidden partition governed by a one-shot pre-registration protocol.
\item
  \textbf{A confirmatory decision-level result inside stated bounds
  (RQ1).} One pre-registered configuration scored once on the sealed
  split: 0.8481 accuracy, 0.8479 macro-F1, Cohen\textquotesingle s κ
  0.7722, with component-cluster intervals roughly 40\% wider than
  item-independence intervals would imply.
\item
  \textbf{Error decomposition by operational harm (RQ3).} The thesis
  separates state-preserving from state-corrupting write failures and
  shows the distinction changes the risk estimate fivefold (0.2111
  aggregate false commit versus 0.0444 state-corrupting), alongside
  missed-update, destructive-overwrite, and unnecessary-abstention rates
  each with finite-sample intervals.
\item
  \textbf{Component attribution (RQ2).} Scoring the rule policy and the
  model policy separately shows neither exceeds 0.750 alone against a
  composed 0.8444, and bounds the abstention guard\textquotesingle s
  hidden-split contribution at most 6.7 accuracy points of the 84.8
  observed.
\item
  \textbf{Computation and cost accounting (RQ3).} Per-decision latency
  distributions, prompt/completion token counts, retry and failure
  counts, and cost-conditional expected loss under three governance
  regimes --- with monetary cost declared unavailable rather than
  reconstructed from list prices.
\item
  \textbf{Measured negative results (RQ4).} The v2
  pilot\textquotesingle s conflict-type shortcut reaches 1.000 held-out
  accuracy; the v2 reference policy does not outperform the structured
  LLM and its preferred status flips with the cost matrix; and the v4.3
  controller\textquotesingle s abstentions are not selectively
  risk-targeted. Each is reported rather than omitted.
\item
  \textbf{Downstream validation with an attributable pattern (RQ1,
  RQ5).} A homogeneous full-slice LongMemEval-M run places governed
  writes at 76.1\% above every structural baseline, with every adjacent
  step in the ordering separating under paired tests and the gain
  concentrated in the competencies the write policy targets.
\item
  \textbf{Evidence exclusion and successor design (RQ4, RQ5).}
  Incomplete, mapped, contaminated, and selection-biased results are
  excluded from primary claims by an explicit ledger; a role-separated
  successor protocol specifies independent generation, three-human
  annotation, long trajectories, hidden baselines scored under the same
  one-shot discipline, and native external metrics.
\item
  \textbf{Reproducible evidence package (RQ5).} Released per-item
  predictions, a recomputation artifact regenerating every reported
  number from raw predictions, audit JSON/Markdown, claim-evidence
  matrix, and machine-readable ledgers keep paper, thesis, code, and
  limitations synchronized.
\end{enumerate}

\subsection{1.6 Scope}\label{scope}

The thesis evaluates the ASK formulation on an author-designed
diagnostic slice and one stratified LongMemEval-M slice with a single
base model family and single seed. It does not claim community-benchmark
status for MGP-6 v2, does not claim end-task superiority on LongMemEval,
does not claim human multi-annotator agreement, and does not claim that
governance on these slices generalizes to open-domain deployments. These
boundaries are stated because they are what the evidence supports;
Chapter 9 enumerates them systematically.

\subsection{1.7 Thesis structure and contribution
map}\label{thesis-structure-and-contribution-map}

{\def\LTcaptype{none} % do not increment counter
\begin{longtable}[]{@{}
  >{\raggedright\arraybackslash}p{(\linewidth - 4\tabcolsep) * \real{0.2030}}
  >{\centering\arraybackslash}p{(\linewidth - 4\tabcolsep) * \real{0.2030}}
  >{\raggedright\arraybackslash}p{(\linewidth - 4\tabcolsep) * \real{0.5940}}@{}}
\toprule\noalign{}
\begin{minipage}[b]{\linewidth}\raggedright
\textbf{Research question}
\end{minipage} & \begin{minipage}[b]{\linewidth}\centering
\textbf{Main chapter(s)}
\end{minipage} & \begin{minipage}[b]{\linewidth}\centering
\textbf{Primary artifact}
\end{minipage} \\
\midrule\noalign{}
\endhead
\bottomrule\noalign{}
\endlastfoot
RQ1 & Chapter 4 & src/ask/ (policy, memory operations, orchestrator) \\
RQ2 & Chapters 5, 7 & benchmarks/mgp6/ (data, manual, metrics), ablation
runs \\
RQ3 & Chapters 6, 7 & metrics\_mgp6.py, ledger known\_metric\_issues \\
RQ4 & Chapters 7, 8, 9 & reports/STATISTICAL\_ANALYSIS.md,
reports/ERROR\_TAXONOMY.md, reports/THREATS\_TO\_VALIDITY.md \\
RQ5 & Chapters 6, 9, App. E--F & artifacts/results\_ledger.json,
reports/CLAIMS\_EVIDENCE\_MATRIX.md \\
\end{longtable}
}

\section{2. Background and Related
Work}\label{background-and-related-work}

This chapter follows the verified related-work audit, with submission
citations rechecked against primary metadata sources on 2026-07-15. The
most important change relative to earlier drafts is a cluster of near
neighbors --- STALE, TOKI, Supersede, MemGuard, and HaluMem --- that are
conceptually closer to ASK than the production systems archived in
Chapter 7 and constrain which novelty claims remain defensible (§2.8).

\subsection{2.1 Memory in language-agent
systems}\label{memory-in-language-agent-systems}

\emph{Memory} covers several mechanisms with different mutability and
deletion semantics: context-window memory, retrieval memory, summary
memory, structured records, and model parameters. Long context does not
remove governance obligations --- long-context models use evidence
unevenly by position (Liu et al., 2024) --- and retrieval-augmented
generation separates storage from generation without providing a write
policy (Lewis et al., 2020; Borgeaud et al., 2022).

The systems line most relevant to ASK treats memory as an agent-managed
store. MemGPT frames the LLM as an operating system that self-edits
tiered memory (Packer et al., 2023); its open-source successor Letta is
evaluated in Chapter 7 (0.678 action accuracy under the favourable
NOOP→ABSTAIN mapping). MemGPT\textquotesingle s memory edits are opaque
function calls: there are no write-time decision criteria, no
abstention, no lineage. Mem0 productionizes an
extract--consolidate--retrieve pipeline whose LLM chooses
ADD/UPDATE/DELETE/NOOP per item (Chhikara et al., 2025); its NOOP means
"nothing new", not "too ambiguous to commit", and it exposes no
supersession lineage (Chapter 7: 0.622 upper bound, invalid-commit rate
1.0). Zep/Graphiti stores a bi-temporal knowledge graph in which
contradicted edges receive invalidation timestamps instead of deletion
(Rasmussen et al., 2025) --- the strongest shipped precedent for
supersession-with-history --- but invalidation is a side effect of LLM
extraction, not a governed decision with an abstain option, and Zep
never scores whether an invalidation was correct. MemOS proposes MemCube
versioning/provenance and uses the word "governance" (Li et al., 2025),
but publishes no decision rule for supersede-vs-keep-vs-refuse; through
its adapter, supersession is not even observable (Chapter 7:
bound-only). A-MEM evolves Zettelkasten-style notes but always
integrates new information (Xu et al., 2025); MemoryBank forgets on an
Ebbinghaus schedule --- time-driven, not evidence-driven (Zhong et al.,
2024); Generative Agents accumulate an append-only stream with
reflection (Park et al., 2023).

Two recent lines are closer to ASK\textquotesingle s decision framing.
Memory-R1 trains an RL memory manager over a discrete
ADD/UPDATE/DELETE/NOOP action space (Yan et al., 2025) --- the clearest
prior formalization of writes as actions --- but the policy is opaque
(PPO/GRPO), the reward is answer-level, and NOOP is not
uncertainty-driven abstention. RMM performs prospective merge-or-add at
session end with retrospective RL-refined retrieval (Tan et al., 2025);
its merge-vs-add echoes supersede-vs-keep but uses similarity thresholds
without contradiction semantics or abstention. MemGuard assigns each
memory an explicit functional role (type) at write time to prevent
cross-type contamination (Ha et al., 2026); ASK\textquotesingle s
no-typing ablation (macro-F1 0.9399→0.5680, KEEP F1→0.0) makes the same
argument experimentally.

\subsection{2.2 Memory lifecycle, revision, and
conflict}\label{memory-lifecycle-revision-and-conflict}

FadeMem classifies new-vs-old memory pairs (compatible / contradictory /
subsumes / subsumed) and applies decay or fusion (Wei et al., 2026); it
always acts on a detected conflict --- there is no abstain --- and
suppresses rather than preserving lineage. PAMU updates preference
signals only when drift is significant, via sliding-window and EMA
fusion (Sun et al., 2025): a conservative, interpretable update trigger,
but for preference vectors only, with no discrete action space. TOKI is
the closest conceptual neighbor found in the verified audit (Wang,
2026): it argues that contradiction resolution \emph{is} write-time
concurrency control, types four production heuristics as bitemporal
operators, and preserves the losing fact in an audit row. Its
await-confirmation operator is a formal ancestor of ABSTAIN and its
audit row parallels supersession lineage. Reddy and
Challaram\textquotesingle s deterministic conflict-resolution result
(Reddy and Challaram, 2026) motivates testing interpretable update logic
on independent data; it does not validate the author-designed MGP-6
pilot by analogy.

\subsection{2.3 Benchmarks and evaluation for agent
memory}\label{benchmarks-and-evaluation-for-agent-memory}

LongMemEval evaluates chat assistants on five long-term memory abilities
including knowledge updates and abstention (Wu et al., 2025); it is this
thesis\textquotesingle s downstream substrate (a stratified 180-item
slice), and it is end-to-end only --- it cannot attribute failure to the
write decision. LoCoMo tests very long conversational memory at scale
(Maharana et al., 2024). HaluMem separately scores extraction, updating,
and QA and reports that hallucinations accumulate at extraction/update
(Chen et al., 2025); it detects hallucination rather than scoring a
governed action choice, and it does not use gold KEEP/SUPERSEDE/ABSTAIN
actions or lineage and premise-resistance metrics. The contribution here
is therefore the operationalization and controlled, layered evaluation
of this particular policy combination, not a claim of primacy for
operation-level evaluation or decision typologies.

STALE is a high-priority benchmark neighbor (Chao et al., 2026): 400
expert-validated implicit-conflict scenarios probing State Resolution,
Premise Resistance, and Implicit Policy Adaptation. It shares
MGP-6\textquotesingle s constructs almost by name
(MGP-6\textquotesingle s query modes include premise\_resistance; its
single annotation disagreement is an implicit\_conflict item). The
decisive difference is scope: STALE probes a frozen
agent\textquotesingle s \emph{read-side} recognition that its memory is
stale, whereas MGP-6 scores \emph{write-time actions} with an abstain
option, supersession lineage, and memory-state metrics. Supersede
isolates the "supersession gap" on the LongMemEval knowledge-update
subset and trains it with GRPO (Patel, 2026); its reward is answer-time
fact currency rather than a write-action taxonomy with abstention,
lineage, or an interpretable policy. GateMem uses the phrase "memory
governance" for a different axis --- access control, deletion
compliance, and utility across principals in shared memory pools (Ren et
al., 2026). RSCB-MC applies risk-sensitive contextual bandits with
explicit abstention to memory \emph{retrieval} in coding agents (Iscan,
2026). Together these verified neighbors preclude any primacy claim for
abstention, governance, update control, provenance, or operation-level
measurement.

Benchmark and dataset documentation follows the spirit of datasheets and
model cards (Gebru et al., 2021; Mitchell et al., 2019); the agreement
statistics of Chapter 5 use Cohen\textquotesingle s κ (Cohen, 1960), and
the interval estimates of Chapters 5--7 use bootstrap methods (Efron and
Tibshirani, 1993).

\subsection{2.4 Belief revision and knowledge
editing}\label{belief-revision-and-knowledge-editing}

The AGM theory of belief change --- expansion, revision, contraction
under minimal-change postulates --- is the formal ancestor of
KEEP/SUPERSEDE (Alchourrón et al., 1985); classical AGM has no native
abstain (non-prioritized revision covers declining input). ASK is best
read as an operational, resource-bounded approximation of
non-prioritized revision for LLM memory; no formal AGM compliance is
claimed or proven. Parametric knowledge editing (ROME, MEMIT) evaluates
edits in weight space with efficacy/generalization/specificity (Meng et
al., 2022; Meng et al., 2023) --- the decision-level-evaluation analogue
for parameters; editing assumes the edit should happen, whereas ASK
governs \emph{whether} to commit at all.

\subsection{2.5 Selective prediction and
abstention}\label{selective-prediction-and-abstention}

Abstention descends from the reject option (Chow, 1970), selective
classification (El-Yaniv and Wiener, 2010; Geifman and El-Yaniv, 2017),
and learning to defer (Madras et al., 2018; Mozannar and Sontag, 2020).
ASK applies that established idea at memory-write time.
TOKI\textquotesingle s await-confirmation and RSCB-MC\textquotesingle s
retrieval-time abstention further narrow the contribution. The
defensible result is a within-generator mechanism signature plus
explicit false-commit/missed-update trade-off, not a general safety
effect.

\subsection{2.6 Personalization and preference
modeling}\label{personalization-and-preference-modeling}

PrefEval shows preference following collapses within about ten turns
zero-shot (Zhao et al., 2025) --- external motivation for the
single-session-preference failure line in Chapter 8, with the difference
that PrefEval measures following while this thesis attributes failure
mechanism (grounding supply vs answer-side discipline). PersonaMem
benchmarks dynamic user profiling across sessions with frontier models
around 50\% (Jiang et al., 2025); LaMP is the canonical
retrieval-augmented personalization suite with static profiles (Salemi
et al., 2023).

\subsection{2.7 Provenance and
auditability}\label{provenance-and-auditability}

Provenance \emph{storage} is well-precedented: Zep\textquotesingle s
bi-temporal validity intervals and episodes (Rasmussen et al., 2025),
TOKI\textquotesingle s audit rows (Wang, 2026), MemOS\textquotesingle s
MemCube metadata (Li et al., 2025), and HaluMem\textquotesingle s
evaluation-side error tracing (Chen et al., 2025). Auditability as a
governance requirement also appears in alignment and risk frameworks
(Ouyang et al., 2022; Bai et al., 2022; Standards and Technology, 2023).
ASK scores provenance and lineage ablations inside its pilot; the
no-lineage result demonstrates that action accuracy can miss lineage
degradation, but author alignment prevents a general causal claim.

\subsection{2.8 Positioning and allowed novelty
claims}\label{positioning-and-allowed-novelty-claims}

Following reports/NOVELTY\_CLAIMS\_ALLOWED.md and the verified audit,
the thesis makes one narrow contribution claim: it
\textbf{operationalizes and evaluates the combination} of an
interpretable KEEP/SUPERSEDE/ABSTAIN write policy, typed and
provenance-tracked records, lineage-preserving transitions, and separate
decision-, stored-state/lineage-, and downstream-answer measurements.
The controlled ablations attribute distinct failure signatures to
components of that combination. The thesis does not claim invention or
primacy for write operations as actions, update/versioning, provenance,
lineage, abstention, write gating, operation-level evaluation, or the
term "memory governance".

\section{3. Problem Formulation}\label{problem-formulation}

\subsection{3.1 System boundary}\label{system-boundary}

Let an interaction stream produce candidate memory items. At step t,
extraction yields a candidate cₜ (typed content with provenance), the
store state is Mₜ (a set of typed records, each active or superseded,
with lineage links), and the controller computes a relation rₜ ∈
\{independent, duplicate, update/contradiction, unclear\} between cₜ and
a selected prior record, then a decision

d\_t = π(c\_t, r\_t, M\_t) ∈ \{KEEP, SUPERSEDE, ABSTAIN\},

with transition semantics

\begin{itemize}
\item
  \textbf{KEEP}: insert cₜ as a new active record (also the correct
  action when no conflicting prior exists);
\item
  \textbf{SUPERSEDE}: insert cₜ, mark the conflicting prior inactive,
  and record a lineage link from the new record to the superseded one;
\item
  \textbf{ABSTAIN}: no mutation; the conflict or ambiguity is recorded
  for review rather than committed.
\end{itemize}

Retrieval reads only active records; lineage remains queryable for audit
and for explicitly historical questions. The policy π may be a
deterministic scorecard (§4.4) or an LLM (§4.5); in both cases the
action vocabulary and the transition semantics are fixed by code, not by
the model.

\subsection{3.2 What makes a write decision
wrong}\label{what-makes-a-write-decision-wrong}

Three governance failure modes are scored separately (formal definitions
in Appendix C):

\begin{enumerate}
\def\labelenumi{\arabic{enumi}.}
\item
  \textbf{Invalid commit} --- committing a write (KEEP or SUPERSEDE) on
  an item whose gold action is ABSTAIN: the store now contains a record
  that policy says should not have been committed.
\item
  \textbf{Missed supersession} --- choosing KEEP or ABSTAIN when the
  gold action is SUPERSEDE: the stale record remains active. Its
  conservative sub-case, SUPERSEDE→ABSTAIN, leaves the store
  consistent-but-stale without committing anything false; the harmful
  sub-case, silent KEEP, leaves an unresolved contradiction active.
\item
  \textbf{Lineage loss} --- performing the right action while failing to
  preserve the supersession chain, invisible to action accuracy but
  fatal to auditability.
\end{enumerate}

Downstream, two state-dependent behaviors are scored: \textbf{stale
retrieval} (a superseded record surfaces as current) and \textbf{premise
resistance} (whether the system rejects a query premise that its
governed memory contradicts).

\subsection{3.3 Functional requirements}\label{functional-requirements}

\begin{itemize}
\item
  \textbf{FR-1 Typed extraction.} Candidate memories carry a type (e.g.,
  stable preference, transient preference, episodic fact),
  holder/subject, and source provenance.
\item
  \textbf{FR-2 Relation classification.} Each candidate is classified
  against retrieved priors before any write.
\item
  \textbf{FR-3 Three-action decision.} Every write decision is one of
  KEEP/SUPERSEDE/ABSTAIN; unknown or unparsable proposals resolve to no
  mutation.
\item
  \textbf{FR-4 Lineage preservation.} SUPERSEDE links the new record to
  the superseded one; superseded records are excluded from ordinary
  retrieval but never silently destroyed.
\item
  \textbf{FR-5 Interpretable policy option.} A deterministic,
  feature-scored policy must be available so decisions can be explained
  and audited without model introspection.
\item
  \textbf{FR-6 Decision export.} Per-item decisions, features, and
  outcomes are exported with stable IDs for scoring.
\item
  \textbf{FR-7 Evidence-labeled reporting.} Every reported metric
  carries the run\textquotesingle s claim-eligibility label.
\end{itemize}

\subsection{3.4 Invariants}\label{invariants}

\begin{enumerate}
\def\labelenumi{\arabic{enumi}.}
\item
  No memory mutation outside the three-action vocabulary.
\item
  No SUPERSEDE without an identified prior record; lineage recorded
  atomically with the write.
\item
  ABSTAIN never mutates state.
\item
  Retrieval of active records never returns superseded ones (violations
  are scored as stale retrieval, not patched).
\item
  No unparsable model output is coerced into a write.
\item
  No result enters a manuscript table without a ledger entry and
  eligibility label.
\end{enumerate}

\subsection{3.5 Threat model and failure
sources}\label{threat-model-and-failure-sources}

Accidental and adversarial sources considered: ambiguous or
contradictory user statements; model hallucination of relations or
targets; distractor content crowding out decisive evidence at retrieval
time; provider failures corrupting extraction (empirically observed:
HTTP-402 no-responses cached as empty extractions in the v7 LongMemEval
run, §7.6); judge errors in downstream grading (empirically observed:
number-word/date false negatives, §8.3); and researcher degrees of
freedom (fix-on-failure loops without confirmation sets --- the reason
the LongMemEval failure replay is excluded from final claims, §9.1).
Chapter 9 maps each to a validity threat with mitigation status.

\subsection{3.6 Failure taxonomy by causal
layer}\label{failure-taxonomy-by-causal-layer}

{\def\LTcaptype{none} % do not increment counter
\begin{longtable}[]{@{}
  >{\raggedright\arraybackslash}p{(\linewidth - 4\tabcolsep) * \real{0.2728}}
  >{\raggedright\arraybackslash}p{(\linewidth - 4\tabcolsep) * \real{0.3636}}
  >{\raggedright\arraybackslash}p{(\linewidth - 4\tabcolsep) * \real{0.3637}}@{}}
\toprule\noalign{}
\begin{minipage}[b]{\linewidth}\raggedright
\textbf{Layer}
\end{minipage} & \begin{minipage}[b]{\linewidth}\centering
\textbf{Example failure}
\end{minipage} & \begin{minipage}[b]{\linewidth}\centering
\textbf{Primary diagnostic}
\end{minipage} \\
\midrule\noalign{}
\endhead
\bottomrule\noalign{}
\endlastfoot
Extraction & candidate fact missed or fabricated & extraction traces;
fail-closed policy (v8) \\
Relation classification & contradiction labeled independent & relation
confusion; escalation gating \\
Decision & KEEP chosen where gold is SUPERSEDE & action confusion
matrix; per-action F1 \\
Commit safety & write on a gold-ABSTAIN item & invalid-commit rate \\
Lineage & supersession chain broken & lineage integrity \\
Retrieval & superseded record surfaces; decisive turn trimmed away &
stale-retrieval rate; SSP taxonomy (§8.1) \\
Answer & correct grounding present, wrong binding & answer-side
categories of the SSP taxonomy \\
Judge & correct answer graded wrong & adjudication ledger; κ between
judges \\
\end{longtable}
}

\section{4. ASK Architecture and
Implementation}\label{ask_architecture_and_implementation}

All code anchors below refer to the canonical repository
(src/ask/\ldots) and were verified in the pipeline reconstruction audit
(docs/audit/CONTROLLERRUN\_RECONSTRUCTION.md).

\subsection{4.1 Architectural overview}\label{architectural-overview}

The controller surrounds a SQLite-backed memory store with a governed
write path: per-turn extraction of typed candidates, relation
classification against retrieved priors, a scored three-action policy,
transition execution with lineage, and retrieval/answer stages that
consume only governed state. The per-session flow is driven by
ControllerOrchestrator.answer() (src/ask/pipeline/orchestrator.py).

\includegraphics[width=6.25in,height=2.67482in,alt={ASK Controller architecture}]{media/image1.png}

ASK Controller architecture

User turns\\
\textbar{}\\
v\\
{[}Extractor (strict JSON schema){]} -\/-\textgreater{} typed candidate
memories\\
\textbar{}\\
v\\
{[}Prefilter: top-64 turns by question similarity{]}\\
\textbar{}\\
v\\
{[}Relation classifier (two-stage, optional escalation){]}\\
\textbar{} candidate vs selected prior\\
v\\
{[}Policy: scored KEEP / SUPERSEDE / ABSTAIN{]}\\
\textbar{}\\
v\\
{[}Memory operations: insert\_active \textbar{} supersede(+lineage)
\textbar{} no-op{]}\\
\textbar{}\\
v\\
{[}Retrieval over active records (dense, k=16, neighbor expansion){]}\\
\textbar{}\\
v\\
{[}Answer-context compilation (session-diverse rerank,
\textasciitilde12k-token budget){]}\\
\textbar{}\\
v\\
{[}Answerer{]} -\/-\textgreater{} {[}Judge / grader with guarded
normalization{]}

\subsection{4.2 Extraction}\label{extraction}

Extractor.render\_prompt() renders a per-turn extraction prompt bound to
a strict JSON schema (ExtractionOutput.model\_json\_schema());
Extractor.extract() is called per user turn, concurrently (default
concurrency 8). Extraction is fail-closed as of v8: on a provider
no-response the extractor raises rather than caching a fabricated empty
extraction --- the v7 behavior of caching empties is the documented
contamination source F1/F2 (§7.6). A content-keyed shared extraction
cache allows cross-run reuse without re-introducing synthetic empties.

\subsection{4.3 Relation classification and prior
selection}\label{relation-classification-and-prior-selection}

For each candidate, \_process\_candidate() runs a two-stage relation
classifier (RelationClassifier.classify()) with an optional cascade
escalation to a stronger model --- gated so that escalation is a no-op
unless a genuinely different model is configured (the v7 run set
BIG\_MODEL == FAST\_MODEL, making its escalation tier inert; disclosed
in §7.6). Prior selection (\_find\_prior, \_select\_prior\_relation)
picks the stored record against which the relation is evaluated.

\subsection{4.4 Prototype scorecard and released reference
policy}\label{prototype-scorecard-and-released-reference-policy}

controller.policy.decide() implements the scored three-action rule:

if S\_abstain \textgreater{} θ\_abs: ABSTAIN\\
elif S\_revise \textgreater{} θ\_sup and relation\_update \textgreater{}
0: SUPERSEDE\\
else: KEEP

with feature scoring in src/ask/controller/scoring.py. The features are
deliberately few and human-legible (a five-feature scorecard over
relation evidence, type compatibility, provenance agreement, and
conflict signals), so a runtime decision can be explained by inspection
of its feature values.

This runtime scorecard did \textbf{not} generate the released MGP-6
reference predictions. Those predictions came from the deterministic
lexical predicate policy in benchmarks/mgp6/systems/policies.py. Chapter
7 therefore evaluates that artifact as a \emph{reference policy}, not as
evidence for the five-feature scorecard or its thresholds. The scorecard
remains an implemented architectural option whose independent tuning and
evaluation are open work.

\subsection{4.5 The LLM policy and the decision
baselines}\label{the-llm-policy-and-the-decision-baselines}

The LLM-backed controller replaces the scorecard with a model decision
(qwen/qwen3.5-9b) inside the same scaffold: same extraction, relation
evidence, action vocabulary, and transition semantics. Two decision
baselines isolate the contribution of the scaffold itself: a zero-shot
LLM decision baseline (no governance framing; 0.700 accuracy,
invalid-commit 0.606) and a six-shot run with 174/180 predictions. The
latter is scored coverage-aware by aligning the six missing IDs as
INVALID (accuracy 0.7833; macro-F1 0.7692) and remains incomplete, not a
completed-run comparison (§7.2).

\subsection{4.6 Memory operations and
lineage}\label{memory-operations-and-lineage}

src/ask/memory/operations.py implements insert\_active, supersede (links
lineage, marks the prior inactive), supersede\_evict, retrieve\_active
(excludes superseded records), and get\_lineage. Lineage is what makes
supersession auditable: the losing record remains queryable with its
link to the superseding record, rather than being deleted or silently
overwritten --- the same auditability argument made formally by
TOKI\textquotesingle s audit rows (Wang, 2026) and operationally by
Zep\textquotesingle s invalidation timestamps (Rasmussen et al., 2025),
here scored directly via the lineage-integrity metric.

\subsection{4.7 Retrieval and answer-context
compilation}\label{retrieval-and-answer-context-compilation}

Retrieval is dense (MiniLM-L6-v2, max sequence length 256, k=16) with
neighbor expansion (±2 turns), assistant/raw-user context pools, and
query-gated historical lineage (historical context is offered only when
the question requests history). Answer-context compilation reranks for
the answer, applies cross-session diversity (v8), and trims to a
\textasciitilde12,000-token budget. The answerer (qwen3.5-9b,
temperature 0) uses per-competency prompts with a WORK→FINAL scaffold.
These retrieval-side mechanisms matter for Chapter 8: grounding supply
is the dominant failure family in the contaminated v7 SSP diagnostic,
while the later top-up shows that results cannot be carried across
configurations.

\subsection{4.8 Judging and
normalization}\label{judging-and-normalization}

LongMemEvalGrader.grade() applies ordered guards (blank, degenerate,
refusal, abstention) before the token guard; \_normalize\_text()
collapses spelled-out cardinals and multiplicatives ("twice"→"2"); the
hard gold-token reject fires only on a provable numeric conflict, with
all lexical/entity/date mismatches deferred to the semantic LLM judge.
These fixes respond to documented judge false negatives (F6, §8.3) and
are themselves part of the audit trail.

\subsection{4.9 Test strategy}\label{test-strategy}

Unit and regression tests cover fail-closed extraction
(test\_extractor\_recovery.py), decision semantics, lineage operations,
merge-by-session-id de-duplication (merge\_run\_shards()), and the guard
logic above. The primary standalone, stdlib-only analysis
(scripts/analysis/audit\_mgp6\_validity.py) recomputes shortcut probes,
held-out metrics, cluster intervals, finite-sample bounds, calibration,
risk--coverage, and cost sensitivity from per-item predictions. The
older mcnemar\_mgp6.py item-level analysis remains reproducible but
secondary (Appendix E).

\subsection{4.10 Operational modes}\label{operational-modes}

The controller supports simulation (no mutation), shadow, assisted
(human approval for selected actions), and autonomous modes; ABSTAIN
routes to review rather than committing. Deployment guidance (Chapter 9)
restricts autonomous supersession to policy-stable, well-audited cases.

\section{5. MGP-6 v2: The Diagnostic
Instrument}\label{mgp-6-v2-the-diagnostic-instrument}

\subsection{5.1 Purpose and status}\label{purpose-and-status}

MGP-6 v2 is an \textbf{author-designed 180-scenario diagnostic slice},
not an independent community benchmark; every claim in this thesis
phrases it as such. Its purpose is to isolate write-time governance
competence --- when to keep, when to supersede, when to abstain ---
under controlled conflict structure, so that decision failures are
attributable and countable. Verified against the raw JSONL on 2026-07-12
(benchmarks/mgp6/data/mgp6\_v2.jsonl).

\subsection{5.2 Composition}\label{composition}

{\def\LTcaptype{none} % do not increment counter
\begin{longtable}[]{@{}
  >{\raggedright\arraybackslash}p{(\linewidth - 2\tabcolsep) * \real{0.3506}}
  >{\centering\arraybackslash}p{(\linewidth - 2\tabcolsep) * \real{0.6494}}@{}}
\toprule\noalign{}
\begin{minipage}[b]{\linewidth}\raggedright
\textbf{Property}
\end{minipage} & \begin{minipage}[b]{\linewidth}\centering
\textbf{Value}
\end{minipage} \\
\midrule\noalign{}
\endhead
\bottomrule\noalign{}
\endlastfoot
Scenarios & 180 \\
Conflict types & T1--T6, 30 each (type definitions:
benchmarks/mgp6/mgp6\_v2\_spec.md; T3 = implicit conflict) \\
Gold actions & SUPERSEDE 90, ABSTAIN 66, KEEP 24 \\
Splits & train 60 / dev 30 / test 90 \\
Query modes & direct\_recall 108, premise\_resistance 45,
implicit\_application 27 \\
\end{longtable}
}

The gold-action distribution is author-imposed rather than
population-derived. SUPERSEDE is the majority class (90/180), while
ABSTAIN and KEEP account for 66 and 24 items. Macro-F1 therefore
accompanies accuracy, but neither metric repairs the more serious
shortcut defect: conflict type deterministically identifies the action
in this pilot. The distribution must not be interpreted as an estimate
of deployment prevalence.

\subsection{5.3 Unit of analysis and split
policy}\label{unit-of-analysis-and-split-policy}

Each scenario has a stable scenario\_id, an interaction context
establishing prior memory, a candidate event, a gold action, a conflict
type, a split, and a query mode for the downstream probe. Evaluation
joins by stable ID. The rule controller was developed against train/dev;
therefore the \textbf{test-split figures (n=90) are always reported
first} and full-corpus figures second (§7.2), a split-hygiene rule
enforced by the claims matrix.

\subsection{5.4 Query modes: probing the state, not just the
decision}\label{query-modes-probing-the-state-not-just-the-decision}

Each scenario carries a downstream probe in one of three modes:
\textbf{direct recall} (does the current value surface?),
\textbf{premise resistance} (the probe asserts a stale or false premise;
a governed system must resist it), and \textbf{implicit application}
(the current value must be used without being named). These modes
connect decision correctness to memory-state behavior; premise
resistance in particular is where the rule and LLM controllers diverge
most sharply (1.0 vs 0.0889, §7.2). The construct is adjacent to
STALE\textquotesingle s read-side premise-resistance probing (Chao et
al., 2026); MGP-6 v2 differs in scoring the \emph{write actions} that
produce the probed state.

\subsection{5.5 Annotation evidence --- stated
precisely}\label{annotation-evidence-stated-precisely}

Two distinct components, never to be conflated (source of truth §1.2):

\begin{enumerate}
\def\labelenumi{\arabic{enumi}.}
\item
  \textbf{One independent annotator} covered the 90-item test split
  (benchmarks/mgp6/annotations/human/): raw agreement \textbf{0.9889}
  {[}0.9667, 1.0{]}, \textbf{Cohen\textquotesingle s κ 0.98146}
  {[}0.93997, 1.0{]}, macro-F1 0.9913, exactly \textbf{one disagreement}
  --- MGP6V2-T3-TEST-016 (implicit conflict; annotator ABSTAIN vs gold
  SUPERSEDE). Caveats: n = 1 annotator; the annotator\textquotesingle s
  independence and human identity rest on author attestation; this is
  \textbf{not} multi-human inter-annotator agreement.
\item
  \textbf{A four-LLM audit panel} (geminiPro, gptSolHigh, kimi, z) over
  the same 90 items: \textbf{Fleiss κ 0.9815}, unanimous on 88/90,
  per-model agreement with gold 1.0 / 0.9778 / 1.0 / 1.0, majority vote
  \textbf{100\%} versus gold. This is a model audit, not human IAA.
\end{enumerate}

Additional disclosures: the human review sheet exists but is 0/180
filled; dataset documentation dated 2026-07-10 predates both annotation
sets (2026-07-11) and wrongly states agreement was unmeasured ---
reconciled in this thesis and flagged for the release. The single
disagreement is retained as an adjudication item and used for
sensitivity discussion (§8.4), not silently relabeled. Because the
labeling manual was authored alongside the scenarios, the panel may
share the author\textquotesingle s reading; this construct-validity
limit is threat C2 in Chapter 9.

\includegraphics[width=6.25in,height=4.77418in,alt={Annotation agreement}]{media/image2.png}

Annotation agreement

\subsection{5.6 Baselines embedded in the
instrument}\label{baselines-embedded-in-the-instrument}

Four reference systems establish what the slice measures: B1 no-memory,
B2 flat retrieval, B3 recency preference, B4 naive overwrite. B2--B4
cannot express abstention (invalid-commit rate 1.0 by construction on
gold-ABSTAIN items), which is precisely the governance gap the
instrument is designed to expose; consequently their comparisons with
the controller are diagnostic, not general-capability results (§7.4,
threat S3).

\subsection{5.7 Relationship to
LongMemEval-M}\label{relationship-to-longmemeval-m}

The stratified LME-M slice (180 items, 6 competencies × 30, seed 0,
built deterministically by scripts/build\_stratified\_slice.py; the
underlying LongMemEval data is external and not redistributed) addresses
the downstream layer: whether governed memory supports later answers
over long histories (Wu et al., 2025). A system can occupy any cell of
the decision-correct × answer-correct table. Chapter 7 uses the archived
downstream runs only as failure analysis because provider contamination,
replay selection, and incomplete coverage prevent a clean estimate of
cross-layer association.

\subsection{5.8 Known metric issues (disclosed, excluded from
claims)}\label{known-metric-issues-disclosed-excluded-from-claims}

As shipped, over\_forgetting\_rate and under\_forgetting\_rate saturate
at 1.0 for every system and are uninformative; they appear in no claim
in this thesis (ledger known\_metric\_issues). The invalid\_commit
metric was revised during the project; all reported numbers use the
revised definition, and pre-revision .bak files are preserved (threat
I6).

\section{6. Experimental Methodology}\label{experimental-methodology}

\subsection{5.9 MGP-6 v4.3: why it was built, how it was built, and what
it can
bear}\label{mgp-6-v4.3-why-it-was-built-how-it-was-built-and-what-it-can-bear}

\subsubsection{5.9.1 The question the instrument has to
answer}\label{the-question-the-instrument-has-to-answer}

Every memory benchmark this thesis surveyed in Chapter 2 scores a
downstream answer. That is the wrong measurement for a write policy,
because a wrong answer can originate in retrieval, in context assembly,
in generation, or in the write decision, and an aggregate answer score
cannot say which. To isolate the write decision the instrument must hold
retrieval, context, and generation \emph{constant} and vary only the
evidential situation. Stated precisely, the question is:

Given the same stored record and the same proposed new value, can a
system

distinguish the case where the proposal is stale (KEEP), the case where
it is

current (SUPERSEDE), and the case where the evidence cannot settle the

question (ABSTAIN)?

That requirement --- not a desire for scale, coverage, or realism ---
dictated the entire construction. MGP-6 v4.3 is small and synthetic
\emph{on purpose}: the control it buys is the only reason it can answer
the question.

\subsubsection{5.9.2 Why v2 had to be
replaced}\label{why-v2-had-to-be-replaced}

The v4.1 shortcut review found that surface lexical and structural
features recovered the intended action without any reasoning over
authority, scope, time, identity, or prior memory. The v2 evidence for
this is decisive and is reported in §7.2: a conflict-type-only lookup
fitted on train+dev reaches \textbf{1.000} held-out accuracy while never
reading the dialogue or the memory state (type-only: 0.9667). In v2,
each conflict family mapped deterministically to one action, so
family-disjoint splits blocked \emph{literal} overlap while leaving the
family→label rule intact across every split. Scenario recognition was
sufficient; evidence reasoning was never tested.

Two lessons carried into the redesign. First, mechanical repairs
(schemas, identifiers, hashes, manifests) had to be kept strictly
separate from the substantive content redesign, so that a passing
structural check could not be mistaken for a repaired construct --- this
separation is what distinguishes v4.1 from v4.2. Second, review had to
happen at the level of the \emph{group of items that share a scenario},
not the individual item: the KEEP, SUPERSEDE, and ABSTAIN variants must
be inspected jointly, because a leak is a property of the contrast
between them, not of any one of them.

\subsubsection{5.9.3 The item}\label{the-item}

Each released record has exactly five fields and the schema sets
additionalProperties: false:

{\def\LTcaptype{none} % do not increment counter
\begin{longtable}[]{@{}
  >{\raggedright\arraybackslash}p{(\linewidth - 2\tabcolsep) * \real{0.3506}}
  >{\raggedright\arraybackslash}p{(\linewidth - 2\tabcolsep) * \real{0.6494}}@{}}
\toprule\noalign{}
\begin{minipage}[b]{\linewidth}\raggedright
\textbf{Field}
\end{minipage} & \begin{minipage}[b]{\linewidth}\centering
\textbf{Content}
\end{minipage} \\
\midrule\noalign{}
\endhead
\bottomrule\noalign{}
\endlastfoot
public\_id & opaque identifier, MGP43- + 20 hex characters \\
memory\_before & the stored record: principal, memory\_slot,
standing\_value \\
dialogue & 9, 11, or 13 turns alternating USER/ASSISTANT, each with a
turn\_id \\
evaluation\_query & fixed instruction naming the slot and the three
permitted actions \\
gold\_action & KEEP, SUPERSEDE, or ABSTAIN \\
\end{longtable}
}

The closed schema is a design commitment, not housekeeping. There is no
conflict-type field, no design target, no component label, no annotator
note --- no channel of the kind that leaked the answer in v2. The
identifier is opaque so that it encodes nothing recoverable.

The dialogue has a fixed rhetorical shape: early turns state the
standing value and name the candidate value; a middle turn declares
which positions in a numbered evidence list are \textbf{indexed} to the
target slot and warns that the remaining positions belong to
neighbouring records; the evidence list is read out; a closing turn
restates the requested value. The decision must be made from the stored
record and this dialogue alone, before any content question is asked.

\subsubsection{5.9.4 The semantic component --- the load-bearing
idea}\label{the-semantic-component-the-load-bearing-idea}

A \textbf{semantic component} is a group of nine items sharing one
principal, one memory\_slot, one standing\_value, and one candidate
value. The nine form a 3 × 3 grid: three dialogue lengths (9, 11, 13
turns) crossed with the three gold actions. Because the grid is crossed,
both margins are exactly balanced. Direct verification on the released
public splits confirms this:

\begin{itemize}
\item
  all \textbf{40/40} public components are exactly 3 KEEP / 3 SUPERSEDE
  / 3 ABSTAIN;
\item
  each dialogue length is exactly \textbf{40 KEEP / 40 SUPERSEDE / 40
  ABSTAIN}.
\end{itemize}

Neither verbosity nor length carries any signal about the action --- a
shortcut that trips many synthetic benchmarks is closed by construction
rather than by audit.

What varies inside a component is narrower still, and this is the
mechanism that makes the instrument work. All nine items draw their
evidence from a \textbf{byte-identical pool of six sentences} ---
verified for all \textbf{120/120} (component, length) triples in the
public data. The three actions differ \emph{only} in \textbf{which two
of those six sentences the dialogue indexes to the target slot}. Across
the 80 triples where the index statement parses mechanically, the three
indexed pairs are always distinct and always \textbf{partition the
six-sentence pool exactly} (80/80).

\subsubsection{5.9.5 A worked component}\label{a-worked-component}

For principal = Rohan, memory\_slot = status message, standing\_value =
available, candidate value do not disturb, all nine items contain the
same six sentences and the same closing turn. Only the index differs:

{\def\LTcaptype{none} % do not increment counter
\begin{longtable}[]{@{}
  >{\raggedright\arraybackslash}p{(\linewidth - 2\tabcolsep) * \real{0.7261}}
  >{\centering\arraybackslash}p{(\linewidth - 2\tabcolsep) * \real{0.2739}}@{}}
\toprule\noalign{}
\begin{minipage}[b]{\linewidth}\raggedright
\textbf{Indexed pair (the two sentences bound to the slot)}
\end{minipage} & \begin{minipage}[b]{\linewidth}\centering
\textbf{Gold}
\end{minipage} \\
\midrule\noalign{}
\endhead
\bottomrule\noalign{}
\endlastfoot
"The do-not-disturb period has an expiry record." + "After the expiry,
Rohan writes that his status is available." & KEEP \\
"Rohan writes that the temporary period has ended." + "He then changes
the status entry to do not disturb." & SUPERSEDE \\
"The retained records do not include the offsets needed to compare
them." + "The expiry time and the status-change time use different time
zones." & ABSTAIN \\
\end{longtable}
}

The reasoning each pair demands differs in kind, which is the point:

\begin{itemize}
\item
  \textbf{KEEP} --- a temporary override expired and the standing value
  was reasserted
\end{itemize}

\emph{after} it. The proposal is stale; the stored record is already
correct.

\begin{itemize}
\item
  \textbf{SUPERSEDE} --- the temporary period ended and a durable change
  followed. The
\end{itemize}

proposal is current; the stored record is out of date.

\begin{itemize}
\item
  \textbf{ABSTAIN} --- the two timestamps are in different time zones
  and the retained
\end{itemize}

records lack the offsets needed to compare them. Neither conclusion is
warranted. Note the character of this case: the evidence is not
\emph{weak} or \emph{sparse}, it is \textbf{unresolvable}. This is
deliberate, and it is why the 300 ABSTAIN items carry a dual ambiguity
flag: they encode designed unresolvability, not annotator uncertainty.

The other four sentences stay in every item as \textbf{controlled
distractors}. A system that reads the whole transcript rather than the
indexed evidence will find apparent support for all three actions in all
three items. The confound is not removed by paraphrasing; it is removed
by holding the surface form fixed and moving only the pointer. That is
the sense in which v4.3 is a \emph{controlled} synthetic benchmark
rather than merely a synthetic one.

\subsubsection{5.9.6 Splits and
determinism}\label{splits-and-determinism}

Components, never items, are assigned to splits: train 180 items / 20
components, development 180 / 20, hidden 540 / 60. Each of the 20
domains contributes one component to train, one to development, and
three to hidden. Assignment uses recorded seed 4301. The build is
byte-deterministic --- archive timestamps are pinned to a constant
(1980-01-01) so identical inputs yield identical hashes. Because a
component never straddles a boundary, no system can observe the KEEP
variant of a scenario in training and the SUPERSEDE variant at test.

\subsubsection{5.9.7 The release gates are executable, not
narrative}\label{the-release-gates-are-executable-not-narrative}

\textbf{Shortcut audit.} Fitted on train, evaluated on hidden,
component-held-out, repeated over five seeds (4301--4305). It trains
word-level and character-level TF-IDF logistic regressions over
\textbf{twelve restricted views} of each item, plus punctuation-only and
length-only numeric features, plus decision trees at depths one through
five. Each view carries a pre-set accuracy ceiling:

{\def\LTcaptype{none} % do not increment counter
\begin{longtable}[]{@{}
  >{\raggedright\arraybackslash}p{(\linewidth - 6\tabcolsep) * \real{0.2831}}
  >{\centering\arraybackslash}p{(\linewidth - 6\tabcolsep) * \real{0.2334}}
  >{\centering\arraybackslash}p{(\linewidth - 6\tabcolsep) * \real{0.2500}}
  >{\centering\arraybackslash}p{(\linewidth - 6\tabcolsep) * \real{0.2334}}@{}}
\toprule\noalign{}
\begin{minipage}[b]{\linewidth}\raggedright
\textbf{View}
\end{minipage} & \begin{minipage}[b]{\linewidth}\centering
\textbf{Ceiling}
\end{minipage} & \begin{minipage}[b]{\linewidth}\centering
\textbf{View}
\end{minipage} & \begin{minipage}[b]{\linewidth}\centering
\textbf{Ceiling}
\end{minipage} \\
\midrule\noalign{}
\endhead
\bottomrule\noalign{}
\endlastfoot
full\_text & 0.65 & assistant\_only & 0.45 \\
dialogue\_only & 0.65 & first\_three & 0.50 \\
target\_evidence\_only & 0.65 & query\_only & 0.40 \\
last\_two & 0.50 & memory\_only & 0.40 \\
candidate\_only & 0.45 & id\_only & 0.40 \\
previous\_turn & 0.45 & punctuation\_only & 0.40 \\
user\_only & 0.45 & length\_only & 0.40 \\
depth3\_tree & 0.60 & & \\
\end{longtable}
}

The release ships with failed\_gates = {[}{]}: no probe cleared its
threshold. Two views deserve specific note. id\_only confirms the opaque
identifiers leak nothing --- the exact class of failure that made
v2\textquotesingle s conflict-type field fatal. target\_evidence\_only
is the strictest self-criticism in the build: it checks that even the
\emph{correctly indexed} evidence, stripped of dialogue context, does
not linearly separate the three classes, so the label depends on reading
the evidence in context rather than on a lexical signature of the
evidence itself.

\textbf{Duplicate audit.} SHA-256 over seven restricted views, failing
on \emph{any} cross-split collision in the five decision-relevant ones
(full\_text, dialogue\_only, candidate\_only, previous\_turn,
last\_two), plus character TF-IDF (4--6 grams) cosine similarity over
every cross-split pair, failing on any edge at ≥ 0.995. Because
components share an evidence pool by construction, within-split
near-duplication is expected and permitted; the gate targets leakage
across the split boundary specifically.

\textbf{Public/private boundary.} Public model-facing records exclude
gold labels, design targets, component metadata, private identifiers,
and annotator metadata. Hidden inputs, hidden gold, final human gold,
completed annotation packets, private mappings, and the detailed
restricted-view audits live only in a gitignored local boundary. Content
is addressed by two Merkle roots --- final gold 437b2a26\ldots{} and
annotations a7e59f88\ldots{} --- and each distributed archive carries
its own SHA-256.

\subsubsection{5.9.8 Human validation and the weight it can
bear}\label{human-validation-and-the-weight-it-can-bear}

Two independent annotation packets cover all 900 items. The builder does
not accept them on trust. Every row must carry a schema version; a
content\_hash matching the coordinator\textquotesingle s binding for
that item; an action in the permitted set; a \textbf{non-empty list of
evidence turn identifiers that all reference real turns in that
dialogue}; a rationale of at least 20 characters; a confidence in
{[}0,1{]}; boolean ambiguity and exclusion flags; an annotator
identifier; and a timezone-aware timestamp. Coverage must be exactly 900
with no duplicate annotation identifiers. Any violation raises
RuntimeError and aborts the build.

The evidence-turn requirement is the substantive one: it makes an
annotation \emph{checkable against the item} rather than a bare label,
and it is what allows the builder to confirm that annotators engaged
with the indexed evidence structure rather than pattern-matching the
scenario.

Outcome: agreement on every action (raw 1.0, Cohen\textquotesingle s κ
1.0), zero exclusions, zero label adjudications, mean confidence 0.96
and 0.967. Both annotators flagged all 300 ABSTAIN items as
intentionally ambiguous scenarios.

\textbf{What this does and does not establish.} Perfect agreement on a
task whose labelling rule is deterministic by construction certifies
that the manual is followable and that the items instantiate it. It does
\textbf{not} independently validate the construct: two annotators
applying the same unambiguous rule will agree whether or not the rule
measures memory governance. This is the honest reading, and it is why κ
= 1.0 appears in this thesis as a construction check rather than as
evidence of difficulty or of external validity. Completion is accepted
through dataset-owner attestation; the builder validates records,
coverage, hashes, and agreement, but does not authenticate reviewer
identity or annotation duration.

\subsubsection{5.9.9 Why the design makes errors
interpretable}\label{why-the-design-makes-errors-interpretable}

Because the three actions within a component are surface-matched, an
error is \emph{diagnostic} rather than merely wrong. Predicting
SUPERSEDE on the KEEP variant means the system read a sentence that was
present in the item but indexed to a different slot --- distractor
capture, specifically.

This is testable, and it was tested. Of the 28 errors the frozen
controller makes on the development split, 18 can be mechanically
resolved against a sibling item\textquotesingle s indexed pair (same
component, same dialogue length, gold action equal to the
\emph{predicted} action). In \textbf{all 18}, the evidence supporting
the predicted action was present in the item but not indexed to the
target slot. Every resolvable error is distractor capture --- precisely
the failure mode the component design was built to isolate. The
instrument is measuring what it was constructed to measure, which is a
stronger statement than any accuracy figure in Chapter 7.

\subsubsection{5.9.10 What the construction cannot
deliver}\label{what-the-construction-cannot-deliver}

The control that makes v4.3 informative is bought by giving up realism,
and the trade must be stated plainly. Deliberately constructed scenarios
do not reproduce the distribution of natural user interactions; the
dialogues are stylized evidence-list recitations, not conversation. The
design is single-slot and single-step, so it cannot test interference
between records or multi-turn write trajectories. It is English-only.
Human completion is owner-attested rather than identity-authenticated.
And the release has no selected public redistribution license, so it
must not be publicly hosted until licensing is resolved. MGP-6 v4.3 is a
controlled diagnostic that must be paired with an external benchmark
such as LongMemEval (§7.11), never a substitute for one.

\subsection{6.1 Evidence policy: the ledger and claim
eligibility}\label{evidence-policy-the-ledger-and-claim-eligibility}

Run provenance comes from artifacts/results\_ledger.json v1.1;
scientific interpretation is superseded by
artifacts/analysis/mgp6\_validity\_audit.json. The failed shortcut gate
makes MGP-6 v2 policy rows diagnostic regardless of their older ledger
label. Mapping-dependent adapters, incomplete runs, and every
LongMemEval row remain non-primary. Annotation is described as one
independent annotator plus a separate model panel, never human IAA.

\subsection{6.2 Systems under test (MGP-6 v2, n =
180)}\label{systems-under-test-mgp-6-v2-n-180}

{\def\LTcaptype{none} % do not increment counter
\begin{longtable}[]{@{}
  >{\raggedright\arraybackslash}p{(\linewidth - 2\tabcolsep) * \real{0.4568}}
  >{\centering\arraybackslash}p{(\linewidth - 2\tabcolsep) * \real{0.5432}}@{}}
\toprule\noalign{}
\begin{minipage}[b]{\linewidth}\raggedright
\textbf{System}
\end{minipage} & \begin{minipage}[b]{\linewidth}\centering
\textbf{Description}
\end{minipage} \\
\midrule\noalign{}
\endhead
\bottomrule\noalign{}
\endlastfoot
B1 no-memory & answers without a store \\
B2 flat retrieval & stores everything, retrieves by similarity; no write
decisions \\
B3 recency & prefers most recent record; no abstention \\
B4 naive overwrite & new information always replaces old; no
abstention \\
Reference policy & deterministic lexical predicate policy used by the
released MGP-6 prediction artifact; not the prototype five-feature
scorecard \\
LLM controller & full ASK with qwen/qwen3.5-9b as the decision policy \\
Zero-shot LLM decision & same model asked for the action without the
governance scaffold \\
Few-shot (6-shot) LLM decision & 174/180 predictions; six missing IDs
scored INVALID for a coverage-aware bound (0.7833 / 0.7692);
incomplete \\
External adapters & Mem0 2.0.11, Letta (MemGPT), Graphiti 0.29.2, MemOS
via adapter harnesses \\
\end{longtable}
}

External-adapter scoring requires an action mapping because these
systems emit NOOP rather than ABSTAIN. Those mapped rankings are
retained in the artifact audit but excluded from primary comparisons. A
successor study must use each system\textquotesingle s native observable
operations and resulting state.

\subsection{6.3 Ablations}\label{ablations}

Four single-component removals from the full rule controller, same items
and scaffold: \textbf{no\_abstain} (forced binary commit),
\textbf{no\_provenance}, \textbf{no\_typing}, \textbf{no\_lineage}.
Their within-generator failure signatures (§7.3) are mechanism
diagnostics: they show how this implementation changes when a switch is
removed, but do not independently establish that each component is
necessary in other datasets or controllers.

\subsection{6.4 LongMemEval-M protocol}\label{longmemeval-m-protocol}

Stratified 180-item slice (6 × 30, seed 0). Five systems (B1--B4,
controller) were run under the v7 configuration (answerer qwen3.5-9b,
judge google/gemma-4-31b-it); the run is \textbf{diagnostic only} due to
documented provider contamination (§7.6). The v8 failure replay
re-executed exactly the 70 v7 controller failures under a changed
configuration and is failure\_replay\_only, not an evaluation of the
full slice. The complementary Modal top-up targeted 110 prior v7
successes but completed only 84, so it is incomplete; it nevertheless
provides decisive counter-evidence to carrying all successes forward
(§7.7). The frozen final-1.1 protocol specifies a homogeneous full-180
rerun with qwen/qwen3.5-9b and a Gemma judge under a fixed provider and
retry policy. Its API smoke failed before inference because the supplied
OpenRouter credential returned HTTP 403, so the run is
externally\_pending and no LongMemEval value is final.

\subsection{6.5 Metrics}\label{metrics}

Decision layer: action accuracy, macro-F1, per-action F1, confusion
counts. Commit safety: invalid-commit rate (share of gold-ABSTAIN items
on which the system commits a write). Memory state: lineage integrity;
stale-retrieval rate. Probe behavior: premise resistance; implicit
application. Downstream (LME-M): judge-graded answer accuracy with
per-competency decomposition and stratified bootstrap CIs. Formal
definitions in Appendix C. Excluded as uninformative: the shipped
over-/under-forgetting rates (§5.8). Premise resistance and
invalid-commit are bespoke metrics defined in metrics\_mgp6.py and are
interpreted descriptively, not as externally validated constructs
(threat C5).

\subsection{6.6 Statistical analysis}\label{statistical-analysis}

The primary analysis uses the held-out 90-item test split and
10,000-replicate paired cluster bootstrap over template\_family\_id,
repeated with minimal-pair groups plus singleton items. Both
item-weighted and family-macro estimates are reported. Exact item-level
McNemar is secondary because within-family independence is doubtful.
Zero event counts receive one-sided finite-sample upper bounds. No
superiority, equivalence, or non-inferiority language is licensed
without a pre-specified contrast and margin. On LongMemEval, no
hypothesis test is computed because no clean comparative design exists.
Reproduction: python3 scripts/analysis/audit\_mgp6\_validity.py.

\subsection{6.7 Leakage and researcher-degrees-of-freedom
controls}\label{leakage-and-researcher-degrees-of-freedom-controls}

Train+development lookup probes over type, conflict type, and public
metadata precede policy interpretation. Conflict type reaches 1.000 test
accuracy and therefore fails the construct-validity gate: test-first
split hygiene alone is insufficient. The v7→v8 fix loop had no
confirmation set, so its output is failure\_replay\_only; overrides
remain separate from raw results; and invalid, incomplete, mapped, and
superseded runs are excluded from primary claims.

\subsection{6.8 MGP-6 v4.3 experimental protocol and completed hidden
run}\label{mgp-6-v4.3-experimental-protocol-and-completed-hidden-run}

MGP-6 v4.3 uses train for loader and engineering checks, and development
for prompt, model, configuration, ablation, cost, and latency selection.
Each system submits one strict record,
\{"public\_id":"MGP43-...","action":"KEEP"\}; the validator rejects
missing, duplicate, unknown, extra, and invalid-action records. Public
loading and prompt construction are deterministic, prediction writing is
resume-safe, and raw-response logs record latency, token, retry, and
failure accounting without private fields.

Before hidden prediction, the release and split hashes, provider and
exact model ID, prompt and system-prompt hashes, decoding, seed, retry
policy, timeout, concurrency, evaluator version, code commit,
environment lock, and prediction schema were frozen and recorded. The
selected public-indexed-evidence ASK live controller with
qwen/qwen3.5-9b was then run exactly once on the 540-item hidden split.
Its prediction file was complete, duplicate-free, and sealed with
SHA-256 e6a598799115c9fc80ae3e731f1b5d59a4947e48df5356b12ba518bafb846806
before the local-private scorer opened hidden gold. A transport
disconnect after 432 durable predictions was recorded; the frozen run
resumed only the remaining 108 records and did not change controls or
duplicate IDs.

The private evaluator reports accuracy, macro-F1, per-class
precision/recall/F1, confusion matrix, coverage, selective and
abstention metrics, and a 5,000-replicate component-cluster bootstrap
with seed 4301. It does not expose gold records in its aggregate output.
Provider telemetry supplies latency, tokens, failures, retries, and
invalid-response counts, but no price field; cost is therefore recorded
as unavailable. LongMemEval remains a separate external evaluation, not
a substitute for this controlled diagnostic. Because only one final v4.3
configuration was pre-registered, there is no final hidden baseline,
ablation, or paired comparison.

\subsubsection{6.8.1 Systems, metrics, and statistical treatment for
v4.3}\label{systems-metrics-and-statistical-treatment-for-v4.3}

The v4.3 comparison covers twelve configurations on the public
development split: the pre-registered frozen controller; the successor
controller in both its as-run and guard-composed forms; four structural
baselines (B1 no memory, B2 flat retrieval, B3 recency, B4 naive
overwrite); an unstructured zero-shot decision policy; the deterministic
public-evidence rule policy in isolation; and four component ablations
(no abstention, no provenance, no typing, no lineage). Three deployed
external memory systems --- Letta, Mem0, and Graphiti --- were run
through adapters written for this thesis. Because their native
operations do not match the KEEP/SUPERSEDE/ABSTAIN vocabulary, their
scores depend on an action mapping and are reported as harness
diagnostics, never as capability measurements of those systems in their
intended use. All model-backed rows use qwen/qwen3.5-9b at temperature
0, seed 4301, 128 maximum new tokens, and concurrency 1.

Primary decision metrics are action accuracy, macro-F1, per-action
precision/recall/F1, Cohen\textquotesingle s κ, and the multiclass
Matthews correlation. The latter two are interpretable here precisely
because every split is exactly action-balanced, so neither is inflated
by a class prior.

Errors are decomposed by \emph{operational harm} rather than reported as
a single rate. Four quantities are tracked, each with its own
denominator: action false commit on gold-ABSTAIN items, separated into
the state-preserving ABSTAIN→KEEP and state-corrupting ABSTAIN→SUPERSEDE
components; missed update on gold-SUPERSEDE items; destructive overwrite
of a gold-KEEP record; and unnecessary abstention over items that
required a decision. Every rate carries a Wilson 95\% interval, and a
zero count is reported with a finite-sample upper bound rather than as
zero risk. Expected loss is computed under balanced, commit-averse, and
update-responsive cost matrices fixed before scoring.

Hidden-split uncertainty uses a 5,000-replicate component-cluster
bootstrap at seed 4301, because items within a semantic component share
a scenario skeleton and are not independent. On the public development
split, component-cluster resampling is \textbf{unavailable} --- the
component map is private and the public scorer does not read it --- so
those intervals are 10,000-replicate item bootstraps and are
consequently \emph{narrower than a component-aware estimate would be}.
Paired system contrasts on development use exact two-sided McNemar tests
with Holm correction over the reported family; they inherit the same
independence assumption. This thesis relies on them only for contrasts
whose effect sizes (b/c splits of 98/6 and larger) could not plausibly
be produced by within-component dependence, and flags the one marginal
contrast --- the 0.117 successor gain --- as the place where the
assumption matters most.

Operational accounting is taken from per-response records rather than
harness wall-clock timing: provider-reported latency per response (mean,
median, p95, maximum), prompt and completion tokens, transport failures,
in-client retries, and invalid-response counts. Monetary cost is
reported only where the provider emitted a price field; on every run in
this thesis it did not, and cost is recorded as unavailable rather than
reconstructed from list-price arithmetic.

\section{7. Results}\label{results}

\subsection{7.1 How to read this
chapter}\label{how-to-read-this-chapter}

This chapter reports two instruments and the reader must not merge them.
\textbf{MGP-6 v4.3} (§7.10) is the successor release whose
construct-validity gate is enforced at build time; its sealed-hidden
result is the thesis\textquotesingle s primary decision-level
measurement, and its development-split comparison is explicitly labelled
development evidence because that split selected the configuration.
\textbf{MGP-6 v2} (§7.2--§7.9) is the earlier pilot whose gate
\emph{failed}; its numbers are retained because the failure is itself a
result and because it motivated the v4.3 redesign, but no v2 row
supports a capability claim. \textbf{LongMemEval-M} (§7.11) is separate
external validation, not a substitute for either.

All values in this chapter are recomputed from released per-item
prediction files rather than transcribed from run reports; the
recomputation is archived in
artifacts/analysis/mgp6\_v43\_full\_benchmark\_analysis.json.

The construct-validity gate precedes every controller result. MGP-6 v2
is an author-designed pilot and its ledger label does not override the
newly measured shortcut failure. Policy scores below are diagnostic
observations on this collection, not estimates of general
memory-governance capability. Mapping-dependent external rows and every
LongMemEval row are excluded from primary comparisons.

\textbf{Shortcut probes fitted on train+development and evaluated on
test:}

{\def\LTcaptype{none} % do not increment counter
\begin{longtable}[]{@{}
  >{\raggedright\arraybackslash}p{(\linewidth - 4\tabcolsep) * \real{0.4589}}
  >{\centering\arraybackslash}p{(\linewidth - 4\tabcolsep) * \real{0.2705}}
  >{\centering\arraybackslash}p{(\linewidth - 4\tabcolsep) * \real{0.2706}}@{}}
\toprule\noalign{}
\begin{minipage}[b]{\linewidth}\raggedright
\textbf{Probe (no dialogue or memory state)}
\end{minipage} & \begin{minipage}[b]{\linewidth}\centering
\textbf{Accuracy}
\end{minipage} & \begin{minipage}[b]{\linewidth}\centering
\textbf{Macro-F1}
\end{minipage} \\
\midrule\noalign{}
\endhead
\bottomrule\noalign{}
\endlastfoot
Type only & 0.9667 & 0.9471 \\
Conflict type only & 1.0000 & 1.0000 \\
\end{longtable}
}

Conflict type deterministically identifies the action. This is the
chapter\textquotesingle s primary result: the pilot cannot distinguish
evidence reasoning from scenario-family recognition. Cluster bootstrap
quantifies dependence inside the collection but cannot repair this
construct failure.

\subsection{7.2 MGP-6 v2 policy
diagnostics}\label{mgp-6-v2-policy-diagnostics}

Test split first (split-hygiene rule), full corpus second.

\textbf{Lexical reference policy.} Test split (n=90): action accuracy
\textbf{0.8889}, macro-F1 \textbf{0.9145}. Full corpus (n=180): 0.9222 /
0.9399. The released predictions were generated by deterministic lexical
predicates, not the prototype five-feature weighted scorecard. On test,
the reference makes 0/33 action false commits, with a one-sided 95\%
upper bound of 0.0868; it misses 10/45 required updates and
unnecessarily abstains on 10/57 non-ABSTAIN items.

\textbf{Structured LLM policy (qwen/qwen3.5-9b).} Test split: 0.9000 /
0.9040. Full corpus: 0.8889 / 0.8913. On test it makes 2/33 action false
commits and misses 7/45 updates. Its test accuracy is numerically above
the reference policy.

\textbf{Secondary full-corpus mechanism table (n = 180; includes train
and development):}

{\def\LTcaptype{none} % do not increment counter
\begin{longtable}[]{@{}
  >{\raggedright\arraybackslash}p{(\linewidth - 8\tabcolsep) * \real{0.2298}}
  >{\centering\arraybackslash}p{(\linewidth - 8\tabcolsep) * \real{0.1591}}
  >{\centering\arraybackslash}p{(\linewidth - 8\tabcolsep) * \real{0.1591}}
  >{\centering\arraybackslash}p{(\linewidth - 8\tabcolsep) * \real{0.2062}}
  >{\centering\arraybackslash}p{(\linewidth - 8\tabcolsep) * \real{0.2458}}@{}}
\toprule\noalign{}
\begin{minipage}[b]{\linewidth}\raggedright
\textbf{System}
\end{minipage} & \begin{minipage}[b]{\linewidth}\centering
\textbf{Accuracy}
\end{minipage} & \begin{minipage}[b]{\linewidth}\centering
\textbf{Macro-F1}
\end{minipage} & \begin{minipage}[b]{\linewidth}\centering
\textbf{Invalid-commit}
\end{minipage} & \begin{minipage}[b]{\linewidth}\centering
\textbf{Eligibility}
\end{minipage} \\
\midrule\noalign{}
\endhead
\bottomrule\noalign{}
\endlastfoot
Lexical reference policy & 0.9222 & 0.9399 & 0.0 & diagnostic \\
Structured LLM policy (qwen3.5-9b) & 0.8889 & 0.8913 & 0.0758 &
diagnostic \\
Zero-shot LLM decision & 0.700 & 0.6657 & 0.606 & diagnostic \\
B4 naive overwrite & 0.4222 & 0.2203 & 1.0 & diagnostic \\
B3 recency & 0.4222 & 0.2203 & 1.0 & diagnostic \\
B1 no-memory & 0.3667 & 0.1789 & --- & diagnostic \\
B2 flat retrieval & 0.1333 & 0.0784 & 1.0 & diagnostic \\
Few-shot (6-shot) decision & 0.7833 & 0.7692 & --- & incomplete: 174/180
predictions; 6 missing scored INVALID \\
\end{longtable}
}

The few-shot row is a \textbf{coverage-aware bound}, not a completed-run
comparison. Its 174 predictions were aligned to the 180 gold IDs and the
six missing IDs were scored INVALID. It is excluded from the
pre-specified significance family and from any claim that assumes
complete coverage.

\subsection{7.3 Ablations (n = 180; within-generator
diagnostics)}\label{ablations-n-180-within-generator-diagnostics}

{\def\LTcaptype{none} % do not increment counter
\begin{longtable}[]{@{}
  >{\raggedright\arraybackslash}p{(\linewidth - 6\tabcolsep) * \real{0.2610}}
  >{\centering\arraybackslash}p{(\linewidth - 6\tabcolsep) * \real{0.2058}}
  >{\centering\arraybackslash}p{(\linewidth - 6\tabcolsep) * \real{0.2058}}
  >{\centering\arraybackslash}p{(\linewidth - 6\tabcolsep) * \real{0.3275}}@{}}
\toprule\noalign{}
\begin{minipage}[b]{\linewidth}\raggedright
\textbf{Configuration}
\end{minipage} & \begin{minipage}[b]{\linewidth}\centering
\textbf{Accuracy}
\end{minipage} & \begin{minipage}[b]{\linewidth}\centering
\textbf{Macro-F1}
\end{minipage} & \begin{minipage}[b]{\linewidth}\centering
\textbf{Distinct failure signature}
\end{minipage} \\
\midrule\noalign{}
\endhead
\bottomrule\noalign{}
\endlastfoot
Full controller & 0.9222 & 0.9399 & --- \\
− abstention & 0.5556 & 0.4278 & invalid-commit 0.0 → 1.0 \\
− provenance & 0.700 & 0.7428 & invalid-commit → 0.606 (identical to the
zero-shot baseline\textquotesingle s rate) \\
− typing & 0.7889 & 0.5680 & KEEP F1 → 0.0 \\
− lineage & 0.9222 & 0.9399 & action labels unchanged; lineage integrity
0.8556 → 0.3689 \\
\end{longtable}
}

The lineage ablation is a useful within-generator mechanism check:
action labels remain identical while the instrument\textquotesingle s
lineage score falls by more than half. It illustrates why decision and
state metrics should be reported separately, but author alignment and
the pilot\textquotesingle s shortcut failure prevent treating the
magnitude as an independent causal estimate.

\subsection{7.4 Held-out clustered comparison and historical item
tests}\label{held-out-clustered-comparison-and-historical-item-tests}

On test, reference minus structured-LLM accuracy is \textbf{−0.0111}.
The paired template-family cluster interval \textbf{{[}−0.0938,
0.0625{]}} includes zero (48 test families); exact item-level McNemar
has b=4, c=5, p=1.0. The reference family-macro interval is {[}0.8125,
0.9688{]} and the structured-LLM interval is {[}0.8333, 0.9688{]}. These
data support neither superiority nor non-inferiority.

The table below preserves the earlier full-corpus item-level analysis
for audit history. It includes train and development items, assumes item
independence, and is not the primary inferential result.

{\def\LTcaptype{none} % do not increment counter
\begin{longtable}[]{@{}
  >{\raggedright\arraybackslash}p{(\linewidth - 12\tabcolsep) * \real{0.1751}}
  >{\centering\arraybackslash}p{(\linewidth - 12\tabcolsep) * \real{0.1369}}
  >{\centering\arraybackslash}p{(\linewidth - 12\tabcolsep) * \real{0.2135}}
  >{\centering\arraybackslash}p{(\linewidth - 12\tabcolsep) * \real{0.1159}}
  >{\centering\arraybackslash}p{(\linewidth - 12\tabcolsep) * \real{0.0985}}
  >{\centering\arraybackslash}p{(\linewidth - 12\tabcolsep) * \real{0.1229}}
  >{\centering\arraybackslash}p{(\linewidth - 12\tabcolsep) * \real{0.1373}}@{}}
\toprule\noalign{}
\begin{minipage}[b]{\linewidth}\raggedright
\textbf{Pair}
\end{minipage} & \begin{minipage}[b]{\linewidth}\centering
\textbf{RD (A−B)}
\end{minipage} & \begin{minipage}[b]{\linewidth}\centering
\textbf{RD 95\% CI (Newcombe)}
\end{minipage} & \begin{minipage}[b]{\linewidth}\centering
\textbf{b / c}
\end{minipage} & \begin{minipage}[b]{\linewidth}\centering
\textbf{OR}
\end{minipage} & \begin{minipage}[b]{\linewidth}\centering
\textbf{Holm p}
\end{minipage} & \begin{minipage}[b]{\linewidth}\centering
\textbf{Decision}
\end{minipage} \\
\midrule\noalign{}
\endhead
\bottomrule\noalign{}
\endlastfoot
Rule vs B4 naive overwrite & +0.5000 & {[}+0.4230, +0.5690{]} & 90 / 0 &
∞ & 4.85e-27 & reject H0 \\
Rule vs B2 flat retrieval & +0.7889 & {[}+0.7181, +0.8382{]} & 142 / 0 &
∞ & 1.44e-42 & reject H0 \\
Rule vs LLM controller & +0.0333 & {[}−0.0188, +0.0873{]} & 14 / 8 &
1.75 & 0.286 & retain H0 (n.s.) \\
LLM controller vs zero-shot & +0.1889 & {[}+0.1096, +0.2664{]} & 46 / 12
& 3.83 & 1.64e-05 & reject H0 \\
\end{longtable}
}

Interpretation, in the order that matters:

\begin{enumerate}
\def\labelenumi{\arabic{enumi}.}
\item
  \textbf{Reference vs structured LLM is a held-out null with the LLM
  numerically higher.} It is not evidence of equality, matching, or
  non-inferiority.
\item
  \textbf{Structured vs zero-shot is descriptive.} The historical
  full-corpus gap is large, but prompt structure and policy scaffolding
  are evaluated on a shortcut-leaking pilot; it is not a clean general
  scaffold-effect estimate.
\item
  \textbf{Reference vs B2/B4 is design-aligned.} The reference encodes
  the labeling manual while the baselines cannot express the required
  actions. These rows are mechanism checks, not capability proofs.
\end{enumerate}

\subsection{7.5 Archived external-system mapping sensitivity
(non-primary)}\label{archived-external-system-mapping-sensitivity-non-primary}

Under the favourable NOOP→ABSTAIN upper-bound mapping (NOOP→KEEP mapping
in parentheses):

{\def\LTcaptype{none} % do not increment counter
\begin{longtable}[]{@{}
  >{\raggedright\arraybackslash}p{(\linewidth - 6\tabcolsep) * \real{0.3040}}
  >{\centering\arraybackslash}p{(\linewidth - 6\tabcolsep) * \real{0.1959}}
  >{\centering\arraybackslash}p{(\linewidth - 6\tabcolsep) * \real{0.1959}}
  >{\centering\arraybackslash}p{(\linewidth - 6\tabcolsep) * \real{0.3042}}@{}}
\toprule\noalign{}
\begin{minipage}[b]{\linewidth}\raggedright
\textbf{System}
\end{minipage} & \begin{minipage}[b]{\linewidth}\centering
\textbf{Accuracy}
\end{minipage} & \begin{minipage}[b]{\linewidth}\centering
\textbf{Macro-F1}
\end{minipage} & \begin{minipage}[b]{\linewidth}\centering
\textbf{Notes}
\end{minipage} \\
\midrule\noalign{}
\endhead
\bottomrule\noalign{}
\endlastfoot
Letta (MemGPT lineage) (Packer et al., 2023) & 0.678 (0.489) & 0.488
(0.302) & \\
Mem0 2.0.11 (Chhikara et al., 2025) & 0.622 (0.400) & 0.442 (0.249) &
invalid-commit 1.0 \\
Graphiti 0.29.2 (Rasmussen et al., 2025) & 0.383 (0.139) & 0.219 (0.104)
& \\
MemOS (Li et al., 2025) & 0.144 & 0.164 & diagnostic\_only: supersession
inexpressible through its adapter \\
\end{longtable}
}

These are harness- and mapping-specific transformations, not fair
primary capability measurements of the systems\textquotesingle{} native
use cases. They are excluded from the paper\textquotesingle s main
comparison. A confirmation study must report native
write/state/retrieval behavior first and show all action mappings only
as sensitivity analyses.

\subsection{7.6 LongMemEval-M: diagnostic results
(diagnostic\_only)}\label{longmemeval-m-diagnostic-results-diagnostic_only}

v7 stratified run (n = 180 each; judge google/gemma-4-31b-it):

{\def\LTcaptype{none} % do not increment counter
\begin{longtable}[]{@{}
  >{\raggedright\arraybackslash}p{(\linewidth - 6\tabcolsep) * \real{0.2941}}
  >{\centering\arraybackslash}p{(\linewidth - 6\tabcolsep) * \real{0.2548}}
  >{\centering\arraybackslash}p{(\linewidth - 6\tabcolsep) * \real{0.2745}}
  >{\centering\arraybackslash}p{(\linewidth - 6\tabcolsep) * \real{0.1766}}@{}}
\toprule\noalign{}
\begin{minipage}[b]{\linewidth}\raggedright
\textbf{System}
\end{minipage} & \begin{minipage}[b]{\linewidth}\centering
\textbf{Correct}
\end{minipage} & \begin{minipage}[b]{\linewidth}\centering
\textbf{Accuracy}
\end{minipage} & \begin{minipage}[b]{\linewidth}\centering
\textbf{SSP}
\end{minipage} \\
\midrule\noalign{}
\endhead
\bottomrule\noalign{}
\endlastfoot
B4 naive overwrite & 114 & 0.6333 & 11/30 \\
ASK controller & 110 & 0.6111 & 11/30 \\
B2 flat vector & 80 & 0.4444 & 6/30 \\
B3 recency & 43 & 0.2389 & 0/30 \\
B1 no-memory & 9 & 0.05 & --- \\
\end{longtable}
}

Two disclosures are mandatory. \textbf{First, the baseline ordering: on
this diagnostic slice, naive overwrite (63.3\%) outperformed the v7
controller (61.1\%).} Any claim of controller superiority on LongMemEval
is unsupported by current evidence. \textbf{Second, the v7 run is
provider-contaminated:} 316 HTTP-402 traces, 388 empty extraction
fallbacks cached as valid, at least 12 manual relaunches, and an
escalation tier that was a no-op (BIG\_MODEL == FAST\_MODEL). Controller
per-competency (v7, post-adjudication 110/180; pre-adjudication
108/180): knowledge-update 21/30, multi-session 15/30,
single-session-assistant 26/30, single-session-preference 11/30,
single-session-user 22/30, temporal-reasoning 15/30. Earlier
partial-range probes (v3: 5/10, 4/10; v4: 17/30; v6: 29/40) are
superseded diagnostics; v5 and the thinking-mode v8 replay are invalid;
both top-up runs are incomplete (§6.7).

\subsection{7.7 Why the failure-replay construction is
excluded}\label{why-the-failure-replay-construction-is-excluded}

The v8 replay re-executed exactly the 70 v7 controller failures under a
changed configuration and judge. It produced 25/70 machine-correct
items. A historical merge then retained all 110 v7 successes, replaced
only the replayed failures, and added four post-hoc Codex/GPT-5
overrides, yielding an archived 139/180 (0.7722) construction. That
number is classified failure\_replay\_only: the selection rule makes
judge and configuration identity perfectly confounded with prior
correctness, the fixes were designed on the same failures, the retained
cases inherited v7 contamination, and the four overrides were not
independent human adjudications. It is not a system result, an estimate
of final-1.1, or a confidence-bound endpoint.

The complementary Modal top-up makes the carry-forward error observable.
It targeted the 110 v7 successes and completed 84: \textbf{71 remained
correct and 13 regressed}. Combining these 84 observations with the 70
replay observations gives \textbf{96/154 raw observed correct}. The 26
unresolved items imply full-180 raw bounds of \textbf{96/180 = 53.3\%}
to \textbf{122/180 = 67.8\%}. If the four replay overrides are shown
separately as an adjudicated sensitivity, the corresponding bounds are
\textbf{55.6\%--70.0\%}. Thus this observed v8 attempt cannot reproduce
the archived carry-forward construction.

For SSP, the complementary evidence is equally important: the top-up
observes \textbf{8/29 correct}, with one unresolved item, so the maximum
is \textbf{9/30}. The v7 diagnostic 11/30 and the
replay\textquotesingle s 0/19 repairs remain useful for causal
investigation, but no SSP count can be carried across configurations.
The homogeneous final-1.1 full-180 run remains externally\_pending: its
smoke request failed before inference with OpenRouter HTTP 403.

\subsection{7.8 Annotation agreement (inferential evidence about the
labels, not the
systems)}\label{annotation-agreement-inferential-evidence-about-the-labels-not-the-systems}

One independent annotator: raw agreement 0.9889 {[}0.9667, 1.0{]},
Cohen\textquotesingle s κ 0.98146 {[}0.93997, 1.0{]}, one disagreement
(MGP6V2-T3-TEST-016). Four-LLM panel: Fleiss κ 0.9815, majority vote
100\% vs gold, 88/90 unanimous. Caveats as in §5.5: n = 1 human + model
panel, not multi-human IAA.

\subsection{7.9 Hypothesis outcomes}\label{hypothesis-outcomes}

H1 is not identified by the shortcut-leaking pilot; H2 supports neither
superiority nor non-inferiority; H3 and H4 remain within-generator
mechanism demonstrations; H5 is unsupported. The independently generated
v3 confirmation protocol is the required gate for stronger conclusions.

\subsection{7.10 MGP-6 v4.3 final hidden result (one frozen
configuration)}\label{mgp-6-v4.3-final-hidden-result-one-frozen-configuration}

\subsubsection{7.10.1 Headline result and its
uncertainty}\label{headline-result-and-its-uncertainty}

The pre-registered ASK live controller attained \textbf{0.8481 accuracy}
and \textbf{0.8479 macro-F1} on all 540 sealed hidden items, coverage
1.000. The component-cluster bootstrap (5,000 resamples, seed 4301)
gives 95\% intervals of \textbf{{[}0.8037, 0.8889{]}} for accuracy and
\textbf{{[}0.8021, 0.8894{]}} for macro-F1.

Because the hidden split is exactly action-balanced (180 KEEP, 180
SUPERSEDE, 180 ABSTAIN), two agreement statistics are interpretable
without prior correction, and both are reported here for the first time:

{\def\LTcaptype{none} % do not increment counter
\begin{longtable}[]{@{}
  >{\raggedright\arraybackslash}p{(\linewidth - 4\tabcolsep) * \real{0.3904}}
  >{\centering\arraybackslash}p{(\linewidth - 4\tabcolsep) * \real{0.2113}}
  >{\raggedright\arraybackslash}p{(\linewidth - 4\tabcolsep) * \real{0.3983}}@{}}
\toprule\noalign{}
\begin{minipage}[b]{\linewidth}\raggedright
\textbf{Quantity}
\end{minipage} & \begin{minipage}[b]{\linewidth}\centering
\textbf{Value}
\end{minipage} & \begin{minipage}[b]{\linewidth}\centering
\textbf{95\% interval}
\end{minipage} \\
\midrule\noalign{}
\endhead
\bottomrule\noalign{}
\endlastfoot
Action accuracy & 0.8481 & {[}0.8037, 0.8889{]} component-cluster \\
--- item-independence reference & & {[}0.8154, 0.8760{]} Wilson \\
Macro-F1 & 0.8479 & {[}0.8021, 0.8894{]} component-cluster \\
Balanced accuracy & 0.8481 & --- \\
Cohen\textquotesingle s κ vs gold & 0.7722 & --- \\
Matthews correlation (multiclass) & 0.7728 & --- \\
Non-abstention coverage & 0.6926 & --- \\
Selective accuracy & 0.8449 & --- \\
\end{longtable}
}

Two observations follow. First, the component-cluster interval is
\textbf{41\% wider} than the Wilson interval that item independence
would imply (width 0.0852 versus 0.0605). That difference is the cost of
the component design, and reporting only the Wilson width would
understate the uncertainty by 29\%. Second, the gap between 0.848
accuracy and 0.772 κ is entirely the \emph{structure} of the errors,
since chance accuracy is exactly 1/3 by construction.

Per-action precision / recall / F1 is ABSTAIN 0.8554 / 0.7889 / 0.8208,
KEEP 0.7947 / 0.8389 / 0.8162, and SUPERSEDE 0.8967 / 0.9167 / 0.9066,
each on support 180. Prediction marginals are 166 ABSTAIN, 190 KEEP, 184
SUPERSEDE.

\subsubsection{7.10.2 Confusion matrix and the error
decomposition}\label{confusion-matrix-and-the-error-decomposition}

{\def\LTcaptype{none} % do not increment counter
\begin{longtable}[]{@{}
  >{\raggedright\arraybackslash}p{(\linewidth - 6\tabcolsep) * \real{0.3530}}
  >{\centering\arraybackslash}p{(\linewidth - 6\tabcolsep) * \real{0.2205}}
  >{\centering\arraybackslash}p{(\linewidth - 6\tabcolsep) * \real{0.1764}}
  >{\centering\arraybackslash}p{(\linewidth - 6\tabcolsep) * \real{0.2501}}@{}}
\toprule\noalign{}
\begin{minipage}[b]{\linewidth}\raggedright
\textbf{Gold \textbackslash{} Predicted}
\end{minipage} & \begin{minipage}[b]{\linewidth}\centering
\textbf{ABSTAIN}
\end{minipage} & \begin{minipage}[b]{\linewidth}\centering
\textbf{KEEP}
\end{minipage} & \begin{minipage}[b]{\linewidth}\centering
\textbf{SUPERSEDE}
\end{minipage} \\
\midrule\noalign{}
\endhead
\bottomrule\noalign{}
\endlastfoot
ABSTAIN & 142 & 30 & 8 \\
KEEP & 18 & 151 & 11 \\
SUPERSEDE & 6 & 9 & 165 \\
\end{longtable}
}

The errors are strongly asymmetric, and for a write-governance system
the asymmetry --- not the total --- is the operational result. SUPERSEDE
is the best-recovered action (0.9167 recall): when the evidence
genuinely establishes a new value, the controller applies it. ABSTAIN is
the worst (0.7889 recall), and its 38 failures split 30 to KEEP and only
8 to SUPERSEDE.

That split is the thesis\textquotesingle s central measurement point
about error taxonomy. An ABSTAIN→KEEP error leaves the stored state
\textbf{byte-identical} to what abstention would have produced; the
memory is not corrupted, and only the audit record that the system saw
ambiguous evidence and declined is lost. An ABSTAIN→SUPERSEDE error
\textbf{actually replaces an active record on insufficient evidence} ---
the precise failure the three-action space exists to prevent. Collapsing
both into a single "false commit" rate of 38/180 = 0.2111 therefore
overstates the harm by roughly fivefold relative to the state-corrupting
component.

{\def\LTcaptype{none} % do not increment counter
\begin{longtable}[]{@{}
  >{\raggedright\arraybackslash}p{(\linewidth - 6\tabcolsep) * \real{0.3005}}
  >{\centering\arraybackslash}p{(\linewidth - 6\tabcolsep) * \real{0.3005}}
  >{\centering\arraybackslash}p{(\linewidth - 6\tabcolsep) * \real{0.1560}}
  >{\centering\arraybackslash}p{(\linewidth - 6\tabcolsep) * \real{0.2429}}@{}}
\toprule\noalign{}
\begin{minipage}[b]{\linewidth}\raggedright
\textbf{Governance error}
\end{minipage} & \begin{minipage}[b]{\linewidth}\centering
\textbf{Count / denominator}
\end{minipage} & \begin{minipage}[b]{\linewidth}\centering
\textbf{Rate}
\end{minipage} & \begin{minipage}[b]{\linewidth}\centering
\textbf{Wilson 95\% CI}
\end{minipage} \\
\midrule\noalign{}
\endhead
\bottomrule\noalign{}
\endlastfoot
Action false commit (gold ABSTAIN acted on) & 38/180 & 0.2111 &
{[}0.1578, 0.2764{]} \\
--- of which ABSTAIN→KEEP (state preserved) & 30/180 & 0.1667 &
{[}0.1193, 0.2280{]} \\
--- of which ABSTAIN→SUPERSEDE (state corrupted) & 8/180 & 0.0444 &
{[}0.0227, 0.0852{]} \\
Missed update (gold SUPERSEDE not applied) & 15/180 & 0.0833 &
{[}0.0511, 0.1329{]} \\
Destructive overwrite (gold KEEP superseded) & 11/180 & 0.0611 &
{[}0.0345, 0.1061{]} \\
Unnecessary abstention (gold KEEP/SUPERSEDE abstained) & 24/360 & 0.0667
& {[}0.0452, 0.0973{]} \\
\end{longtable}
}

Unnecessary abstention at 0.0667 is low, which matters: the controller
is not purchasing its conservatism by refusing to act. But selective
accuracy (0.8449) is \emph{below} overall accuracy (0.8481), so the
items it does commit on are answered slightly less accurately than
average. A well-calibrated selective predictor shows the opposite
ordering. \textbf{The abstention mechanism is therefore doing governance
work --- avoiding unjustified writes --- without doing selective-risk
work}, and this thesis does not conflate the two. This is the clearest
single defect in the frozen configuration and it is disclosed rather
than omitted.

\subsubsection{7.10.3 Component decomposition: which part does the
work}\label{component-decomposition-which-part-does-the-work}

The controller composes a public-evidence rule policy with a live model
policy under a conservative guard that overrides a model ABSTAIN only
when the rule policy finds decisive public evidence. Each component is
separately scorable on development data.

{\def\LTcaptype{none} % do not increment counter
\begin{longtable}[]{@{}
  >{\raggedright\arraybackslash}p{(\linewidth - 4\tabcolsep) * \real{0.3568}}
  >{\centering\arraybackslash}p{(\linewidth - 4\tabcolsep) * \real{0.3451}}
  >{\centering\arraybackslash}p{(\linewidth - 4\tabcolsep) * \real{0.2981}}@{}}
\toprule\noalign{}
\begin{minipage}[b]{\linewidth}\raggedright
\end{minipage} & \begin{minipage}[b]{\linewidth}\centering
\textbf{Public dev (n=180)}
\end{minipage} & \begin{minipage}[b]{\linewidth}\centering
\textbf{Hidden (n=540)}
\end{minipage} \\
\midrule\noalign{}
\endhead
\bottomrule\noalign{}
\endlastfoot
Model action alone & 0.7500 & not scorable† \\
Rule policy alone & 0.7500 & not scorable† \\
Composed (reported) & 0.8444 & 0.8481 \\
Items where the guard changed the action & 27 & 36 \\
--- repaired a model error & 20 & --- \\
--- broke a correct model action & 3 & --- \\
--- error-neutral & 4 & --- \\
\end{longtable}
}

† Component-wise hidden scoring would require a second pass over sealed
gold beyond the single authorised scoring event.

On the hidden split the guard fired 36 times and always out of ABSTAIN
(30 → KEEP, 6 → SUPERSEDE). This \textbf{bounds} its contribution: even
under the most favourable accounting, the guard accounts for at most 6.7
accuracy points of the 84.8 observed. The rules that fired were
no-decisive-public-evidence (255), direct-documented-slot-change (120),
standing-value-confirmed-or-proposal- retracted (90),
unresolved-scope-authority-or-timing (57), and
temporary-or-scoped-exception (18). The result is a genuine composition,
not a rule engine with a model attached nor a model with cosmetic rules.

\subsubsection{7.10.4 Cost sensitivity across governance
regimes}\label{cost-sensitivity-across-governance-regimes}

{\def\LTcaptype{none} % do not increment counter
\begin{longtable}[]{@{}
  >{\raggedright\arraybackslash}p{(\linewidth - 6\tabcolsep) * \real{0.2779}}
  >{\centering\arraybackslash}p{(\linewidth - 6\tabcolsep) * \real{0.1953}}
  >{\centering\arraybackslash}p{(\linewidth - 6\tabcolsep) * \real{0.2439}}
  >{\centering\arraybackslash}p{(\linewidth - 6\tabcolsep) * \real{0.2829}}@{}}
\toprule\noalign{}
\begin{minipage}[b]{\linewidth}\raggedright
\textbf{System}
\end{minipage} & \begin{minipage}[b]{\linewidth}\centering
\textbf{Balanced}
\end{minipage} & \begin{minipage}[b]{\linewidth}\centering
\textbf{Commit-averse}
\end{minipage} & \begin{minipage}[b]{\linewidth}\centering
\textbf{Update-responsive}
\end{minipage} \\
\midrule\noalign{}
\endhead
\bottomrule\noalign{}
\endlastfoot
Frozen controller, hidden (n=540) & 0.1519 & 0.4204 & 0.2963 \\
Frozen controller, dev & 0.1556 & 0.4167 & 0.3000 \\
Successor controller, dev & 0.0500 & 0.1500 & 0.1333 \\
Successor + composed guard, dev & 0.0389 & 0.0889 & 0.0944 \\
B4 naive overwrite, dev & 0.6611 & 3.2333 & 1.6222 \\
Rule policy alone, dev & 0.6667 & 1.2444 & 1.6667 \\
\end{longtable}
}

The commit-averse loss is the largest for the frozen controller because
that regime charges most heavily for exactly its dominant error class.
Critically, \textbf{the ordering among ASK controllers is stable across
all three regimes} --- unlike the v2 reference-versus-LLM comparison
(§7.4), where the preferred system \emph{flips} with the cost matrix.
This is a stronger result than the v2 chapter could produce and it is
stated as such.

\subsubsection{7.10.5 Computation time and token
cost}\label{computation-time-and-token-cost}

All runs used concurrency 1, temperature 0, and a 128-token generation
cap, so these are per-item serial costs and not a throughput claim.

{\def\LTcaptype{none} % do not increment counter
\begin{longtable}[]{@{}
  >{\raggedright\arraybackslash}p{(\linewidth - 14\tabcolsep) * \real{0.1575}}
  >{\centering\arraybackslash}p{(\linewidth - 14\tabcolsep) * \real{0.0931}}
  >{\centering\arraybackslash}p{(\linewidth - 14\tabcolsep) * \real{0.1240}}
  >{\centering\arraybackslash}p{(\linewidth - 14\tabcolsep) * \real{0.1097}}
  >{\centering\arraybackslash}p{(\linewidth - 14\tabcolsep) * \real{0.0954}}
  >{\centering\arraybackslash}p{(\linewidth - 14\tabcolsep) * \real{0.0954}}
  >{\centering\arraybackslash}p{(\linewidth - 14\tabcolsep) * \real{0.1527}}
  >{\centering\arraybackslash}p{(\linewidth - 14\tabcolsep) * \real{0.1722}}@{}}
\toprule\noalign{}
\begin{minipage}[b]{\linewidth}\raggedright
\textbf{Run}
\end{minipage} & \begin{minipage}[b]{\linewidth}\centering
\textbf{n}
\end{minipage} & \begin{minipage}[b]{\linewidth}\centering
\textbf{mean (ms)}
\end{minipage} & \begin{minipage}[b]{\linewidth}\centering
\textbf{median}
\end{minipage} & \begin{minipage}[b]{\linewidth}\centering
\textbf{p95}
\end{minipage} & \begin{minipage}[b]{\linewidth}\centering
\textbf{max}
\end{minipage} & \begin{minipage}[b]{\linewidth}\centering
\textbf{prompt tok/item}
\end{minipage} & \begin{minipage}[b]{\linewidth}\centering
\textbf{completion tok/item}
\end{minipage} \\
\midrule\noalign{}
\endhead
\bottomrule\noalign{}
\endlastfoot
Frozen controller, hidden & 540 & 1,673 & 1,278 & 2,931 & 12,626 & 393 &
7.4 \\
Frozen controller, dev & 180 & 1,400 & 959 & 3,669 & 10,260 & 396 &
7.5 \\
Successor controller, dev & 180 & 1,288 & 952 & 3,388 & 7,593 & 506 &
7.8 \\
Zero-shot decision, dev & 180 & 1,425 & 1,039 & 3,122 & 9,886 & 549 &
8.0 \\
B4 naive overwrite, dev & 180 & 1,647 & 1,139 & 4,116 & 23,586 & 578 &
7.9 \\
B2 flat retrieval, dev & 180 & 1,581 & 1,060 & 4,582 & 8,478 & 597 &
7.9 \\
B3 recency, dev & 180 & 1,774 & 1,071 & 5,515 & 11,404 & 233 & 7.9 \\
Rule policy (no model call), dev & 180 & 0.24 & 0.13 & 0.64 & 3.6 & 0 &
0 \\
\end{longtable}
}

The write decision is cheap. Hidden-split totals are 212,467 prompt
tokens, 4,000 completion tokens, and 903.5 s of summed model latency
across 540 items. The latency distribution is right-skewed with a
\emph{thin} tail rather than a fat one: median-to-p95 is a factor of 2.3
while p95-to-max is 4.3, so the worst cases are route-level stalls, not
hard items. These are provider-reported wall times through a shared
OpenRouter route and are \textbf{not} compute measurements; the
deterministic rule policy, which issues no network call, completes the
same 180 decisions at a 0.13 ms median.

Token cost separates the systems more informatively than latency.
Evidence-indexed routing spends 393 prompt tokens per hidden decision
against 578 for naive overwrite, 597 for flat retrieval, and 549 for
zero-shot --- all of which score at chance. The successor configuration
costs 28\% more than the frozen one (506 vs 396 prompt tokens on dev)
and buys 10.6 accuracy points.

\textbf{Monetary cost is unavailable, not estimated.} OpenRouter
telemetry carried no price or cost field on any run in this thesis.
Given the frozen model identifier a list-price multiplication would be
arithmetically trivial and epistemically empty: routing, provider
selection, and caching all affect the billed amount and none of that
appears in the logs. Token counts are recorded because they were
measured.

Run integrity on the hidden split: 0 in-client retries, 0 invalid or
unparseable responses, 0 duplicate predictions, and one transport-level
disconnect after 432 predictions recovered by the immutable-manifest
resume under identical controls. One caveat on the manifest --- its
207.96 s wall\_clock\_seconds field records only the resumed invocation
that produced the final 108 predictions, so it is \textbf{not} the
end-to-end duration of the 540-item run. The 903.5 s summed model
latency is the defensible aggregate.

\subsubsection{7.10.6 Development-split system comparison (n = 180,
balanced)}\label{development-split-system-comparison-n-180-balanced}

Because only one hidden configuration was authorised, every system
contrast is development evidence on the split used for selection. Chance
accuracy is 0.3333.

{\def\LTcaptype{none} % do not increment counter
\begin{longtable}[]{@{}
  >{\raggedright\arraybackslash}p{(\linewidth - 14\tabcolsep) * \real{0.2162}}
  >{\centering\arraybackslash}p{(\linewidth - 14\tabcolsep) * \real{0.1104}}
  >{\centering\arraybackslash}p{(\linewidth - 14\tabcolsep) * \real{0.1323}}
  >{\centering\arraybackslash}p{(\linewidth - 14\tabcolsep) * \real{0.1395}}
  >{\centering\arraybackslash}p{(\linewidth - 14\tabcolsep) * \real{0.1031}}
  >{\centering\arraybackslash}p{(\linewidth - 14\tabcolsep) * \real{0.0994}}
  >{\centering\arraybackslash}p{(\linewidth - 14\tabcolsep) * \real{0.0994}}
  >{\centering\arraybackslash}p{(\linewidth - 14\tabcolsep) * \real{0.0998}}@{}}
\toprule\noalign{}
\begin{minipage}[b]{\linewidth}\raggedright
\textbf{System}
\end{minipage} & \begin{minipage}[b]{\linewidth}\centering
\textbf{Acc.}
\end{minipage} & \begin{minipage}[b]{\linewidth}\centering
\textbf{95\% CI}
\end{minipage} & \begin{minipage}[b]{\linewidth}\centering
\textbf{Macro-F1}
\end{minipage} & \begin{minipage}[b]{\linewidth}\centering
\textbf{MCC}
\end{minipage} & \begin{minipage}[b]{\linewidth}\centering
\textbf{FC}
\end{minipage} & \begin{minipage}[b]{\linewidth}\centering
\textbf{MU}
\end{minipage} & \begin{minipage}[b]{\linewidth}\centering
\textbf{DO}
\end{minipage} \\
\midrule\noalign{}
\endhead
\bottomrule\noalign{}
\endlastfoot
Successor controller (guard-composed) & 0.9611 & {[}0.933, 0.989{]} &
0.9607 & 0.943 & 0.000 & 0.017 & 0.033 \\
Successor controller (as run) & 0.9500 & {[}0.917, 0.978{]} & 0.9494 &
0.926 & 0.000 & 0.033 & 0.067 \\
Frozen controller (pre-registered) & 0.8444 & {[}0.789, 0.894{]} &
0.8458 & 0.771 & 0.233 & 0.133 & 0.033 \\
B1 no memory & 0.3667 & {[}0.294, 0.439{]} & 0.3118 & 0.065 & 0.800 &
0.217 & 0.733 \\
B2 flat retrieval & 0.3556 & {[}0.283, 0.428{]} & 0.2440 & 0.065 & 0.900
& 0.067 & 0.917 \\
B3 recency & 0.3278 & {[}0.261, 0.400{]} & 0.2175 & −0.015 & 0.933 &
0.117 & 0.900 \\
B4 naive overwrite & 0.3389 & {[}0.272, 0.411{]} & 0.1891 & 0.023 &
1.000 & 0.017 & 0.967 \\
Zero-shot LLM decision & 0.3389 & {[}0.272, 0.411{]} & 0.1788 & 0.046 &
1.000 & 0.000 & 0.983 \\
Rule policy alone (no model) & 0.3333 & {[}0.267, 0.406{]} & 0.2434 &
0.000 & 0.217 & 0.783 & 0.217 \\
− abstention & 0.3333 & {[}0.267, 0.406{]} & 0.2465 & 0.000 & 1.000 &
0.233 & 0.767 \\
− provenance & 0.3389 & {[}0.272, 0.411{]} & 0.1887 & 0.029 & 0.983 &
0.017 & 0.967 \\
− typing & 0.3444 & {[}0.278, 0.417{]} & 0.1899 & 0.092 & 0.983 & 0.000
& 0.983 \\
− lineage & 0.3500 & {[}0.283, 0.422{]} & 0.2020 & 0.088 & 0.983 & 0.000
& 0.967 \\
Letta (adapter-mapped) & 0.2944 & {[}0.228, 0.361{]} & 0.2424 & −0.069 &
0.417 & 0.733 & 0.383 \\
Mem0 (adapter-mapped) & 0.3333 & {[}0.267, 0.406{]} & 0.1851 & 0.000 &
0.033 & 1.000 & 0.000 \\
Graphiti (adapter-mapped) & 0.3056 & {[}0.239, 0.372{]} & 0.2112 &
−0.064 & 0.217 & 0.867 & 0.117 \\
\end{longtable}
}

FC = action false commit on gold-ABSTAIN; MU = missed update on
gold-SUPERSEDE; DO = destructive overwrite of a gold-KEEP record.

The most instructive fact here is negative. \textbf{Every} baseline,
ablation, zero-shot policy, standalone rule policy, and deployed
external system lands between 0.2944 and 0.3667 --- at or below chance
--- with macro-F1 in 0.179--0.312 and MCC in −0.069 to 0.092. They are
close to uninformative on this task, and they fail in two
\emph{opposite} ways. Commit-happy systems (B2, B4, zero-shot, and the
no-abstention/no-typing/no-lineage ablations) supersede almost
everything: destructive-overwrite rates 0.767--0.983, false-commit rates
0.900--1.000. Abstain-happy systems (rule policy alone, Mem0, Graphiti)
do the reverse: Mem0 abstains on 117/120 items requiring a decision and
misses all 60 required updates. Both score near chance while behaving in
opposite ways --- which is precisely why this thesis reports the error
decomposition rather than accuracy alone.

Paired exact McNemar tests, Holm-corrected over the 15-test family:

{\def\LTcaptype{none} % do not increment counter
\begin{longtable}[]{@{}
  >{\raggedright\arraybackslash}p{(\linewidth - 12\tabcolsep) * \real{0.1966}}
  >{\centering\arraybackslash}p{(\linewidth - 12\tabcolsep) * \real{0.1715}}
  >{\centering\arraybackslash}p{(\linewidth - 12\tabcolsep) * \real{0.1062}}
  >{\centering\arraybackslash}p{(\linewidth - 12\tabcolsep) * \real{0.1062}}
  >{\centering\arraybackslash}p{(\linewidth - 12\tabcolsep) * \real{0.1162}}
  >{\centering\arraybackslash}p{(\linewidth - 12\tabcolsep) * \real{0.1413}}
  >{\centering\arraybackslash}p{(\linewidth - 12\tabcolsep) * \real{0.1620}}@{}}
\toprule\noalign{}
\begin{minipage}[b]{\linewidth}\raggedright
\textbf{A}
\end{minipage} & \begin{minipage}[b]{\linewidth}\centering
\textbf{B}
\end{minipage} & \begin{minipage}[b]{\linewidth}\centering
\textbf{b}
\end{minipage} & \begin{minipage}[b]{\linewidth}\centering
\textbf{c}
\end{minipage} & \begin{minipage}[b]{\linewidth}\centering
\textbf{RD}
\end{minipage} & \begin{minipage}[b]{\linewidth}\centering
\textbf{Holm p}
\end{minipage} & \begin{minipage}[b]{\linewidth}\centering
\textbf{Decision}
\end{minipage} \\
\midrule\noalign{}
\endhead
\bottomrule\noalign{}
\endlastfoot
Frozen controller & Letta & 106 & 7 & +0.5500 & 1.2e−22 & reject H0 \\
Frozen controller & − abstention & 98 & 6 & +0.5111 & 2.2e−21 & reject
H0 \\
Frozen controller & B1 no memory & 90 & 4 & +0.4778 & 4.2e−21 & reject
H0 \\
Frozen controller & B4 naive overwrite & 98 & 7 & +0.5056 & 1.4e−20 &
reject H0 \\
Frozen controller & Rule policy alone & 106 & 14 & +0.5111 & 5.8e−18 &
reject H0 \\
Frozen controller & Mem0 & 106 & 14 & +0.5111 & 5.8e−18 & reject H0 \\
Successor (guard-comp.) & Frozen controller & 28 & 7 & +0.1167 & 1.5e−03
& reject H0 \\
Successor (as run) & Frozen controller & 28 & 9 & +0.1056 & 5.1e−03 &
reject H0 \\
Successor (guard-comp.) & Successor (as run) & 2 & 0 & +0.0111 & 0.500 &
retain H0 \\
\end{longtable}
}

(Remaining contrasts against B2, B3, zero-shot, − provenance, − typing,
− lineage, and Graphiti all reject at Holm p \textless{} 1.6e−19; full
table in the released analysis artifact.)

The guard-composition repair therefore did \textbf{not} measurably
improve the successor system; it restored the intended composition
semantics without damaging accuracy, and this thesis claims no more than
that.

\subsubsection{7.10.7 What this result is
not}\label{what-this-result-is-not}

This is one measurement of one configuration on a controlled synthetic
diagnostic after public-development model selection. It is \textbf{not}
a comparison against the development-only baselines or ablations under
sealed conditions, not a paired hidden test, and not evidence of
real-world memory-governance performance or deployment safety. The
successor configuration\textquotesingle s 0.9611 development accuracy
has not been tested on any sealed split and must not be read as an
improvement over the 0.8481 hidden result: they are different splits
under different exposure regimes. LongMemEval (§7.11) remains a separate
external-validation track.

\subsubsection{7.11 LongMemEval-M V12: the homogeneous full-slice
run}\label{longmemeval-m-v12-the-homogeneous-full-slice-run}

The V12 run (freeze final-2.0) evaluates all 180 stratified items with
one answer model (qwen/qwen3.5-9b), one judge (google/gemma-4-31b-it),
and zero judge-model fallbacks. It is the eligible downstream result on
this track.

{\def\LTcaptype{none} % do not increment counter
\begin{longtable}[]{@{}
  >{\raggedright\arraybackslash}p{(\linewidth - 10\tabcolsep) * \real{0.1424}}
  >{\centering\arraybackslash}p{(\linewidth - 10\tabcolsep) * \real{0.1392}}
  >{\centering\arraybackslash}p{(\linewidth - 10\tabcolsep) * \real{0.1908}}
  >{\centering\arraybackslash}p{(\linewidth - 10\tabcolsep) * \real{0.1779}}
  >{\centering\arraybackslash}p{(\linewidth - 10\tabcolsep) * \real{0.1844}}
  >{\centering\arraybackslash}p{(\linewidth - 10\tabcolsep) * \real{0.1653}}@{}}
\toprule\noalign{}
\begin{minipage}[b]{\linewidth}\raggedright
\textbf{System}
\end{minipage} & \begin{minipage}[b]{\linewidth}\centering
\textbf{Correct/n}
\end{minipage} & \begin{minipage}[b]{\linewidth}\centering
\textbf{Accuracy (95\% CI)}
\end{minipage} & \begin{minipage}[b]{\linewidth}\centering
\textbf{Prompt tok/item}
\end{minipage} & \begin{minipage}[b]{\linewidth}\centering
\textbf{Latency/item (s)}
\end{minipage} & \begin{minipage}[b]{\linewidth}\centering
\textbf{vs. row above}
\end{minipage} \\
\midrule\noalign{}
\endhead
\bottomrule\noalign{}
\endlastfoot
B1 no memory & 9/180 & 0.0500 {[}0.022, 0.083{]} & 95 & 1.6 & --- \\
B3 recency & 43/180 & 0.2389 {[}0.183, 0.300{]} & 768 & 3.1 & 4.7e−10 \\
B2 flat vector & 80/180 & 0.4444 {[}0.372, 0.517{]} & 3,629 & 2.5 &
3.7e−06 \\
B4 naive overwrite & 114/180 & 0.6333 {[}0.561, 0.700{]} & 14,227 &
966.0 & 6.8e−06 \\
ASK governed writes & 137/180 & 0.7611 {[}0.700, 0.822{]} & 22,625 &
809.2 & 2.9e−04 \\
\end{longtable}
}

Every adjacent step in the monotone ordering separates under an exact
paired McNemar test, Holm-corrected over the six tests reported. The
controller-over- naive-overwrite contrast is b = 31, c = 8, RD = +0.128,
Holm p = 2.9e−04.

Per competency (30 items each):

{\def\LTcaptype{none} % do not increment counter
\begin{longtable}[]{@{}
  >{\raggedright\arraybackslash}p{(\linewidth - 10\tabcolsep) * \real{0.3325}}
  >{\centering\arraybackslash}p{(\linewidth - 10\tabcolsep) * \real{0.1314}}
  >{\centering\arraybackslash}p{(\linewidth - 10\tabcolsep) * \real{0.1314}}
  >{\centering\arraybackslash}p{(\linewidth - 10\tabcolsep) * \real{0.1314}}
  >{\centering\arraybackslash}p{(\linewidth - 10\tabcolsep) * \real{0.1314}}
  >{\centering\arraybackslash}p{(\linewidth - 10\tabcolsep) * \real{0.1419}}@{}}
\toprule\noalign{}
\begin{minipage}[b]{\linewidth}\raggedright
\textbf{Competency}
\end{minipage} & \begin{minipage}[b]{\linewidth}\centering
\textbf{B1}
\end{minipage} & \begin{minipage}[b]{\linewidth}\centering
\textbf{B3}
\end{minipage} & \begin{minipage}[b]{\linewidth}\centering
\textbf{B2}
\end{minipage} & \begin{minipage}[b]{\linewidth}\centering
\textbf{B4}
\end{minipage} & \begin{minipage}[b]{\linewidth}\centering
\textbf{ASK}
\end{minipage} \\
\midrule\noalign{}
\endhead
\bottomrule\noalign{}
\endlastfoot
knowledge update & 0.133 & 0.667 & 0.567 & 0.700 & 0.967 \\
multi-session & 0.033 & 0.200 & 0.133 & 0.567 & 0.633 \\
single-session assistant & 0.000 & 0.000 & 0.900 & 0.867 & 0.967 \\
single-session preference & 0.000 & 0.000 & 0.200 & 0.367 & 0.533 \\
single-session user & 0.067 & 0.400 & 0.767 & 0.767 & 0.833 \\
temporal reasoning & 0.067 & 0.167 & 0.100 & 0.533 & 0.633 \\
\end{longtable}
}

The gain over naive overwrite is \textbf{not uniform}, and its shape is
the strongest available evidence that the improvement comes from the
write decision rather than from generic pipeline differences. It
concentrates in knowledge update (+8 items) and single-session
preference (+5) --- exactly where a stale value must be replaced or a
durable preference separated from a transient one. On multi-session the
gain is +2 items and on temporal reasoning +3, both within noise at this
slice size. A better retriever would not preferentially help
knowledge-update items.

\textbf{Write-decision behaviour.} The controller\textquotesingle s
action log records 12,654 write decisions across 3,547 sessions: 11,478
KEEP (90.7\%), 779 ABSTAIN (6.2\%), and 397 SUPERSEDE (3.1\%) over
candidate types episodic\_fact (9,857), transient\_pref (2,416), and
stable\_pref (381). This distribution is what the design predicts on
natural dialogue, where most turns carry no conflicting evidence, and it
differs sharply from the balanced MGP-6 v4.3 diagnostic --- which is the
reason both instruments are retained. All 397 supersessions recorded
both a predecessor and a successor identifier, and none of the 779
abstentions produced a resulting record, so the lineage and no-op
invariants held on all 12,654 decisions.

\subsubsection{Pipeline anatomy: where the compute actually
goes}\label{pipeline-anatomy-where-the-compute-actually-goes}

The released pipeline traces (extraction\_trace.jsonl,
relation\_trace.jsonl, answers.jsonl) permit a stage-level decomposition
rather than a single per-item cost figure, and the decomposition
overturns the natural reading of the accuracy table.

{\def\LTcaptype{none} % do not increment counter
\begin{longtable}[]{@{}
  >{\raggedright\arraybackslash}p{(\linewidth - 10\tabcolsep) * \real{0.1582}}
  >{\centering\arraybackslash}p{(\linewidth - 10\tabcolsep) * \real{0.1549}}
  >{\centering\arraybackslash}p{(\linewidth - 10\tabcolsep) * \real{0.1549}}
  >{\centering\arraybackslash}p{(\linewidth - 10\tabcolsep) * \real{0.1817}}
  >{\centering\arraybackslash}p{(\linewidth - 10\tabcolsep) * \real{0.1817}}
  >{\centering\arraybackslash}p{(\linewidth - 10\tabcolsep) * \real{0.1685}}@{}}
\toprule\noalign{}
\begin{minipage}[b]{\linewidth}\raggedright
\textbf{Stage}
\end{minipage} & \begin{minipage}[b]{\linewidth}\centering
\textbf{Model calls}
\end{minipage} & \begin{minipage}[b]{\linewidth}\centering
\textbf{Prompt tok.}
\end{minipage} & \begin{minipage}[b]{\linewidth}\centering
\textbf{Completion tok.}
\end{minipage} & \begin{minipage}[b]{\linewidth}\centering
\textbf{Share of prompt}
\end{minipage} & \begin{minipage}[b]{\linewidth}\centering
\textbf{Model latency}
\end{minipage} \\
\midrule\noalign{}
\endhead
\bottomrule\noalign{}
\endlastfoot
Extraction & 11,847 & 41,806,968 & 1,806,685 & 73.6\% & 19.98 h \\
Relation classification & 19,157 & 10,936,290 & 1,226,921 & 19.2\% &
18.93 h \\
Answer generation & 180 & 4,072,477 & 41,418 & 7.2\% & see note \\
Total & 31,184 & 56,815,735 & 3,075,024 & --- & 38.91 h \\
Per item & 173.2 & 315,643 & 17,083 & --- & 778 s \\
\end{longtable}
}

Answering the 180 questions takes 180 model calls and 4.07 M prompt
tokens. Reaching the point where those questions \emph{can} be answered
takes 31,004 further calls and 52.7 M prompt tokens. The write-time
stages are \textbf{92.8\% of prompt tokens and 99.4\% of model calls};
the answer stage --- the only stage an end-to-end memory benchmark
scores --- is 7.2\% of the budget.

This reframes the cost comparison above. The governed controller is not
an expensive answerer; it is a cheap answerer sitting on an expensive
ingestion pipeline, and the same holds for naive overwrite. Reporting
only answer-stage tokens, as the retrieval-only baselines invite,
measures the 7\% and ignores the 93\%.

\textbf{A correction to how per-item latency must be read.} The
latency\_ms field on each answer record is \textbf{cumulative
end-to-end} time for that item, not answer-generation time. Summing it
gives 40.46 h, while summed extraction and relation model latency is
38.91 h; the 1.55 h difference is the implied generation-only cost ---
about \textbf{31 s per item}, not the 809 s the field suggests. Per-item
pipeline times quoted in this section use the cumulative field and must
be read as end-to-end pipeline time. This thesis flags the point
explicitly because the naive reading --- that one answer call takes
thirteen minutes --- is wrong by more than an order of magnitude.

\textbf{Extraction is mostly rejection.} Of 11,491 processed turns,
\textbf{5,656 (49.2\%)} return no\_extractable\_content. The
extractor\textquotesingle s most common output is \emph{nothing}, which
is correct behaviour on conversational filler and the first line of
defence against memory bloat. The remaining turns yield 11,771 candidate
memories (1.02 per turn overall), typed as 9,144 episodic facts, 2,261
transient preferences, and 366 stable preferences. Mean extraction
confidence is 0.934 (minimum 0.75); 17.4\% of candidates carry a hedge
marker and 1.3\% a negation marker. Parsing was robust: 370 of 11,491
turns hit a first-parse error, 9 required the repair path, and there
were \textbf{zero model fallbacks and zero retries beyond the first
attempt}.

\textbf{Relation routing is deliberately conservative --- and
auditable.} The stage is two-part: stage one labels the candidate
against a retrieved prior as support, contradict, or independent; stage
two, invoked on 9,530 of 9,539 classifications, resolves same-subject
and temporal ordering. Routing is then deterministic and every decision
records its reason:

{\def\LTcaptype{none} % do not increment counter
\begin{longtable}[]{@{}
  >{\raggedright\arraybackslash}p{(\linewidth - 4\tabcolsep) * \real{0.3890}}
  >{\centering\arraybackslash}p{(\linewidth - 4\tabcolsep) * \real{0.4283}}
  >{\centering\arraybackslash}p{(\linewidth - 4\tabcolsep) * \real{0.1827}}@{}}
\toprule\noalign{}
\begin{minipage}[b]{\linewidth}\raggedright
\textbf{Stage-1 label}
\end{minipage} & \begin{minipage}[b]{\linewidth}\centering
\textbf{Routed relation}
\end{minipage} & \begin{minipage}[b]{\linewidth}\centering
\textbf{n}
\end{minipage} \\
\midrule\noalign{}
\endhead
\bottomrule\noalign{}
\endlastfoot
independent & unclear & 8,281 \\
support & support & 880 \\
contradict & supersede & 221 \\
contradict & contradict & 147 \\
contradict & unclear & 1 \\
parse fallback & unclear & 9 \\
\end{longtable}
}

All 8,281 independent labels route to unclear --- an independent
candidate is never permitted to displace a prior. Of 369 contradict
labels, only the 221 that stage two confirms as \textbf{same-subject and
newer} route to supersede; 113 older and 33 same-time contradictions are
held as contradict without a write. So \textbf{221 of 9,539
classifications (2.3\%)} may propose a replacement. This is the
mechanism behind the 397 supersessions (3.1\%) in the action log and,
ultimately, behind the controller\textquotesingle s low
destructive-overwrite rate --- the conservatism is architectural, not
emergent.

\textbf{Caching did not help, and the trace explains why.} The
extraction cache holds 11,477 unique keys against 11,491 traced turns,
the relation cache 9,523 against 9,539 classifications --- under 0.3\%
deduplication in both. Each turn is keyed by its own content and each
candidate by its pairing with a specific prior, so in a single pass over
distinct sessions there is essentially nothing to reuse. The caches are
resume insurance, not a throughput optimisation, and any deployment
estimate assuming cache amortisation across a first pass is unsupported
by this evidence.

\textbf{Downstream cost in context.} Per-item end-to-end pipeline time
is 809 s for the governed controller and 966 s for naive overwrite,
against 1.6--3.1 s for the three retrieval-only baselines --- roughly
320× against flat-vector retrieval, buying 31.7 accuracy points.
Answer-stage prompt tokens show the same ordering (22,625 per item
versus 14,227 and 3,629), but as the stage decomposition above shows,
that comparison understates the true gap by more than an order of
magnitude, because the retrieval-only baselines have no ingestion
pipeline to account for. Whether the exchange is worthwhile is a
deployment question this thesis does not adjudicate; the full accounting
is reported so it is not hidden behind the accuracy ordering.

\textbf{Judge accounting.} Every item was scored by the single frozen
judge: 40 items resolved by deterministic guards (normalised match, no
LLM call), 140 routed to the LLM judge, 2 scoring fallbacks logged, 0
model fallbacks. Nine items received a documented semantic-equivalence
adjudication correcting judge scoring errors, each with a written
rationale and a preserved pre-adjudication backup.

\textbf{Two open provenance items} do not affect any reported number:
the homogeneity gate flags a stale code/protocol hash for the evaluation
launcher, and the manifest was not written to a clean terminal status
after a resume. Per-system integrity checks (answer model, judge
routing, fallback counts) pass with empty error and warning lists.

\section{8. Error Analysis and
Discussion}\label{error-analysis-and-discussion}

\subsection{8.1 Single-session preference (SSP): a v7 diagnostic
taxonomy}\label{single-session-preference-ssp-a-v7-diagnostic-taxonomy}

In the contaminated v7 run, SSP is the controller\textquotesingle s
weakest competency: \textbf{11/30 correct (0.3667), CI {[}0.20,
0.5333{]}}. The failure-only replay repaired 0/19 of those v7 failures,
but it used different prompt contexts and therefore does not isolate an
answer-model effect. The complementary top-up observes \textbf{8/29
correct} with one unresolved item (maximum \textbf{9/30}), proving that
the v7 count cannot be carried forward as a cross-version result. The
taxonomy below is consequently a causal analysis of the \textbf{19 v7
failures}, not a final-1.1 evaluation
(docs/research/SINGLE\_SESSION\_PREFERENCE\_FAILURE\_ANALYSIS.md).

\textbf{Primary failure categories (n = 19):}

{\def\LTcaptype{none} % do not increment counter
\begin{longtable}[]{@{}
  >{\raggedright\arraybackslash}p{(\linewidth - 4\tabcolsep) * \real{0.5262}}
  >{\centering\arraybackslash}p{(\linewidth - 4\tabcolsep) * \real{0.2369}}
  >{\centering\arraybackslash}p{(\linewidth - 4\tabcolsep) * \real{0.2370}}@{}}
\toprule\noalign{}
\begin{minipage}[b]{\linewidth}\raggedright
\textbf{Primary category}
\end{minipage} & \begin{minipage}[b]{\linewidth}\centering
\textbf{Count}
\end{minipage} & \begin{minipage}[b]{\linewidth}\centering
\textbf{Share}
\end{minipage} \\
\midrule\noalign{}
\endhead
\bottomrule\noalign{}
\endlastfoot
retrieval-omission & 6 & 31.6\% \\
cross-session-contamination & 4 & 21.1\% \\
target-attribute-extraction-miss & 3 & 15.8\% \\
distractor-entity-binding & 2 & 10.5\% \\
polarity-error & 2 & 10.5\% \\
unsupported-answer-specificity & 1 & 5.3\% \\
immediate-request-vs-enduring-preference & 1 & 5.3\% \\
\end{longtable}
}

\textbf{Combined mechanism mentions (primary + secondary, of 19):}
retrieval-omission 11; cross-session-contamination 8;
distractor-entity-binding 7; target-attribute-extraction-miss 5;
unsupported-answer-specificity 3;
immediate-request-vs-enduring-preference 2; polarity-error 2.

\textbf{Pipeline-layer localization:} retrieval-ranking 12/19 (63.2\%);
answer-realization 10/19 (52.6\%); context-construction (top-k trim)
5/19 (26.3\%); attribution/target-binding 5/19 (26.3\%); ingestion,
extraction, persistence, controller-policy, and judge-scorer 0/19 each.
Notably, \textbf{no SSP failure was traced to the write-governance
policy itself}: the decisive facts exist in the corpus, and for the
answer-side cases they were even present in the prompt.

Two families dominate. \textbf{Grounding-supply failures}
(retrieval-omission + cross-session-contamination: primary in 10/19,
present in 15/19): the single supporting session\textquotesingle s
decisive preference turn is out-ranked by topically similar distractor
sessions and trimmed away, and the model then answers from the wrong
session or fabricates. The 0/19 replay repair is consistent with that
localization but is not causal proof because the replay changed prompt
contexts as well as the answer path. \textbf{Answer-side
binding/realization failures} (primary in 9/19): the decisive fact was
in context, but the answer bound the preference to the wrong target,
violated an explicitly stated negative constraint (recommending a
dis-preferred podcast genre), or over-specified beyond evidence. One
item (1da05512) is a borderline judge-severity case rather than a clean
pipeline defect, and several judge rationales say "fabricates" where the
mechanism is distractor binding --- the offending tokens are verbatim in
retrieved memory.

\textbf{Highest-leverage proposed fix (hypothesized, not run):} make
preference grounding an explicit, provenance-constrained representation
--- a typed record (holder / target / attribute / polarity / scope /
source-session) whose retrieval is guaranteed for the
question\textquotesingle s subject --- paired with an answer contract
that asserts only retrieved-and-grounded preferences and honors stated
negatives. The design targets the \textasciitilde15/19 grounding-supply
mentions, while the \textasciitilde9/19 answer-side failures show why a
retrieval-only repair is unlikely to suffice. Necessity and effect size
remain untested. The design and its credit-attribution ablation matrix
are specified in docs/research/SINGLE\_SESSION\_PREFERENCE\_METHOD.md;
PrefEval\textquotesingle s finding that preference following collapses
within ten turns (Zhao et al., 2025) independently motivates the same
construct.

\subsection{8.2 MGP-6 error profiles: three systems, three
signatures}\label{mgp-6-error-profiles-three-systems-three-signatures}

\textbf{Reference policy --- conservative, one-directional on the full
pilot.} Its 14 full-corpus action errors are SUPERSEDE→ABSTAIN. On test
it has 0/33 action false commits, but the one-sided 95\% upper bound is
0.0868 and it misses 10/45 updates. This is a measured trade-off inside
the pilot, not a population-safety result.

\textbf{Structured LLM policy --- different error allocation.} Test
action accuracy is 0.9000 versus 0.8889 for the reference, with a paired
family interval containing zero. It has more action false commits (2/33)
but fewer missed updates (7/45). This supports cost-sensitive reporting,
not an absolute policy ordering.

\textbf{Zero-shot decision baseline --- diagnostic gap.} Its lower score
and high full-corpus invalid-commit rate show a failure mode worth
studying, but label leakage prevents attributing the difference cleanly
to the governance scaffold.

\textbf{Ablation signatures.} Abstention removal → invalid-commit 1.0;
typing removal → KEEP F1 0.0; lineage removal → audit-score decay
invisible to action accuracy (0.8556→0.3689). These are distinct
within-generator failure signatures, consistent with the motivation
behind MemGuard\textquotesingle s typing argument (Ha et al., 2026) and
TOKI\textquotesingle s audit rows (Wang, 2026). They motivate
independent component tests; they do not by themselves establish
necessity, causal effect size, or external validity.

\subsection{8.3 LongMemEval run-integrity failure classes
(F1--F12)}\label{longmemeval-run-integrity-failure-classes-f1f12}

The failure ledger (docs/audit/CONTROLLERRUN\_FAILURE\_LEDGER.md)
distinguishes per-item failures (§8.1) from run/pipeline failures, of
which the most consequential are: \textbf{F1/F2} --- HTTP-402 provider
no-responses cached as valid empty extractions (316 "402" traces; 388
empty extractions), fixed in v8 by fail-closed extraction; \textbf{F9}
--- relation-classifier no-responses converted to a cached semantic
Keep, same fix; \textbf{F3/F4} --- total-loss replay failures
(thinking-mode token exhaustion → 0 predictions; embedder segfaults),
fixed by enable\_thinking:false and a shared embedding server;
\textbf{F6} --- judge false negatives on number-words and dates ("4" vs
"four", "Valentine\textquotesingle s Day" vs "February 14th"), fixed by
guarded normalization with deferral to the semantic judge; \textbf{F8}
--- the v7 escalation no-op; \textbf{F5/F11} --- provenance defects
(misnamed aborted run; a manifest pointing to a non-existent replay
directory); \textbf{F10} --- in-place verdict overrides preserved with
originals; and the top-up evidence showing that prior successes cannot
be retained by assumption. The general lesson: infrastructure failures
that default to "empty but valid" silently corrupt memory state, while
selective replay silently biases reported accuracy; fallbacks must fail
closed and repairs must be evaluated on homogeneous coverage.

\subsection{8.4 The retained annotation disagreement as construct
evidence}\label{the-retained-annotation-disagreement-as-construct-evidence}

MGP6V2-T3-TEST-016 (implicit conflict): gold SUPERSEDE, independent
annotator ABSTAIN. The disagreement sits on the same SUPERSEDE/ABSTAIN
boundary as the reference policy\textquotesingle s conservative errors,
making it a useful candidate for independent adjudication. With one
human label and an author-aligned manual, it is not enough to identify a
generally hard construct region. It is retained for blinded
adjudication; sensitivity is bounded (one item = 1.1\% of the test
split) and both labels are defensible under the current manual.

\subsection{8.5 Decision--state--answer
divergence}\label{decisionstateanswer-divergence}

The pilot illustrates why the three layers should be reported
separately. Within MGP-6, similar action accuracy can coexist with very
different premise-resistance behavior, and the no-lineage ablation
leaves action labels unchanged while changing the
instrument\textquotesingle s state metric. Separately, the contaminated
LongMemEval slice happens to rank B4 naive overwrite above the
controller on answers (63.3\% vs 61.1\%). Because the first observations
are author-aligned and the second is contaminated, these are diagnostic
examples rather than clean cross-layer effect estimates. They
nevertheless show how a single aggregate can conceal different failure
locations and motivate the layered evaluation design (RQ3).

\subsection{8.6 Asymmetric harm and the price of
conservatism}\label{asymmetric-harm-and-the-price-of-conservatism}

The reference policy\textquotesingle s full-pilot action errors are
conservative, but missed supersession can itself be harmful. Whether the
trade is right is policy-dependent: under the frozen commit-averse
matrix the reference has lower test loss (0.111 vs 0.189), while under
the update-responsive matrix the structured LLM has lower loss (0.433 vs
0.556). The current weights and any prototype score thresholds have not
been validated on independent data, so this is sensitivity analysis
rather than optimization evidence.

\subsection{8.7 What the near-neighbor cluster
implies}\label{what-the-near-neighbor-cluster-implies}

Five verified 2026-adjacent works (STALE, Supersede, TOKI, MemGuard,
HaluMem) cover read-side staleness probing, the supersession gap, formal
write-time contradiction operators, write-time typing, and
operation-level evaluation. This thesis is positioned as an
\textbf{evaluation decomposition and validity audit}: it makes no claim
that the individual mechanisms or their conjunction are unprecedented.
The convergence supports the importance of the problem, while the
deterministic-policy finding of (Reddy and Challaram, 2026) motivates
independent confirmation rather than validating this pilot by analogy.

\section{9. Validity, Ethics, and
Limitations}\label{validity-ethics-and-limitations}

This chapter follows the v1.1 evidence policy and the repository threat
analysis; every threat cites a concrete artifact and carries a
mitigation status.

\subsection{9.1 Internal validity}\label{internal-validity}

{\def\LTcaptype{none} % do not increment counter
\begin{longtable}[]{@{}
  >{\raggedright\arraybackslash}p{(\linewidth - 6\tabcolsep) * \real{0.1397}}
  >{\raggedright\arraybackslash}p{(\linewidth - 6\tabcolsep) * \real{0.2866}}
  >{\raggedright\arraybackslash}p{(\linewidth - 6\tabcolsep) * \real{0.2866}}
  >{\raggedright\arraybackslash}p{(\linewidth - 6\tabcolsep) * \real{0.2870}}@{}}
\toprule\noalign{}
\begin{minipage}[b]{\linewidth}\raggedright
\textbf{\#}
\end{minipage} & \begin{minipage}[b]{\linewidth}\centering
\textbf{Threat}
\end{minipage} & \begin{minipage}[b]{\linewidth}\centering
\textbf{Instance}
\end{minipage} & \begin{minipage}[b]{\linewidth}\centering
\textbf{Status}
\end{minipage} \\
\midrule\noalign{}
\endhead
\bottomrule\noalign{}
\endlastfoot
I1 & Dataset/controller circularity and label leakage & conflict type
alone reaches 1.000 held-out accuracy & failed construct gate; not
repaired by policy comparisons or cluster bootstrap; independent v3 data
required \\
I2 & Failure-derived fixes without a confirmation set (LongMemEval
v7→v8) & fixes designed on the exact 70 failures then replayed on them &
mitigated for reporting: replay classified failure\_replay\_only; open
experimentally: homogeneous rerun required \\
I3 & Adjudication contamination & v7: 108 original judge-correct + 2
manual corrections; replay merge: 4 Codex/GPT-5 overrides that are not
independent human adjudication & mitigated for reporting: raw/overridden
counts separated; overrides excluded from final claims \\
I4 & Provider contamination and carry-forward bias & v7 has 316 HTTP-402
traces and 388 empty fallbacks; top-up finds 13 regressions among 84
observed prior successes & mitigated for reporting: no v7/v8 result
claim; open experimentally: homogeneous rerun \\
I5 & Split hygiene (MGP-6) & controller developed on train/dev;
type/action mapping persists across splits & not sufficient: test
figures are primary but remain diagnostic \\
I6 & Metric revision mid-project & invalid\_commit semantics revised;
.bak preserved & mitigated: revised definition used throughout \\
\end{longtable}
}

\subsection{9.2 External validity}\label{external-validity}

{\def\LTcaptype{none} % do not increment counter
\begin{longtable}[]{@{}
  >{\raggedright\arraybackslash}p{(\linewidth - 4\tabcolsep) * \real{0.1705}}
  >{\raggedright\arraybackslash}p{(\linewidth - 4\tabcolsep) * \real{0.4146}}
  >{\raggedright\arraybackslash}p{(\linewidth - 4\tabcolsep) * \real{0.4149}}@{}}
\toprule\noalign{}
\begin{minipage}[b]{\linewidth}\raggedright
\textbf{\#}
\end{minipage} & \begin{minipage}[b]{\linewidth}\centering
\textbf{Threat}
\end{minipage} & \begin{minipage}[b]{\linewidth}\centering
\textbf{Status}
\end{minipage} \\
\midrule\noalign{}
\endhead
\bottomrule\noalign{}
\endlastfoot
E1 & MGP-6 v2 is an author-designed diagnostic slice, not a community
benchmark & open: mapped external-system comparisons do not provide
independent validation; v3 protocol required \\
E2 & LongMemEval evaluated on a 180-item slice, not the full 500 &
open \\
E3 & Single base model (qwen3.5-9b) and single seed; CIs reflect item
sampling, not run variance & open \\
E4 & External-adapter numbers are harness- and mapping-specific &
partially mitigated: both NOOP mappings reported; MemOS flagged
bound-only \\
E5 & Synthetic-template dependence (templates, labels, and reference
policy authored together) & confirmed threat: conflict/type shortcut
succeeds despite family-disjoint splits \\
\end{longtable}
}

\subsection{9.3 Construct validity}\label{construct-validity}

{\def\LTcaptype{none} % do not increment counter
\begin{longtable}[]{@{}
  >{\raggedright\arraybackslash}p{(\linewidth - 4\tabcolsep) * \real{0.1705}}
  >{\raggedright\arraybackslash}p{(\linewidth - 4\tabcolsep) * \real{0.4146}}
  >{\raggedright\arraybackslash}p{(\linewidth - 4\tabcolsep) * \real{0.4149}}@{}}
\toprule\noalign{}
\begin{minipage}[b]{\linewidth}\raggedright
\textbf{\#}
\end{minipage} & \begin{minipage}[b]{\linewidth}\centering
\textbf{Threat}
\end{minipage} & \begin{minipage}[b]{\linewidth}\centering
\textbf{Status}
\end{minipage} \\
\midrule\noalign{}
\endhead
\bottomrule\noalign{}
\endlastfoot
C1 & "Human IAA" would be a mis-description: evidence is n=1 human +
4-LLM panel & partially mitigated by the mandated phrasing; open:
multi-human IAA; review sheet 0/180 \\
C2 & Panel may share the author\textquotesingle s reading (manual
unambiguous by construction) & open \\
C3 & Shipped forgetting metrics uninformative (saturate at 1.0) &
mitigated by disclosure: excluded from claims \\
C4 & Judge as construct proxy on LongMemEval (known false-negative
classes) & partially mitigated: guard fixes, deferral to semantic judge,
logged adjudication \\
C5 & Premise resistance and invalid-commit are bespoke metrics &
partially mitigated: defined in code; interpreted descriptively \\
\end{longtable}
}

\subsection{9.4 Statistical-conclusion
validity}\label{statistical-conclusion-validity}

{\def\LTcaptype{none} % do not increment counter
\begin{longtable}[]{@{}
  >{\raggedright\arraybackslash}p{(\linewidth - 4\tabcolsep) * \real{0.1705}}
  >{\raggedright\arraybackslash}p{(\linewidth - 4\tabcolsep) * \real{0.4146}}
  >{\raggedright\arraybackslash}p{(\linewidth - 4\tabcolsep) * \real{0.4149}}@{}}
\toprule\noalign{}
\begin{minipage}[b]{\linewidth}\raggedright
\textbf{\#}
\end{minipage} & \begin{minipage}[b]{\linewidth}\centering
\textbf{Threat}
\end{minipage} & \begin{minipage}[b]{\linewidth}\centering
\textbf{Status}
\end{minipage} \\
\midrule\noalign{}
\endhead
\bottomrule\noalign{}
\endlastfoot
S1 & Small and clustered n (90 test items, 48 test families; SSP strata
30) & partially mitigated: family/minimal-pair bootstrap and exact
counts; external uncertainty remains open \\
S2 & Multiplicity & mitigated: Holm over the pre-specified 4-test
family \\
S3 & Tautological rule-vs-weak-baseline comparisons (OR = ∞ by design) &
mitigated by framing: "not sampling noise", not capability proof \\
S4 & The reference-vs-LLM null is easy to misreport as equality or
non-inferiority & mitigated: test RD −0.0111 and family interval
reported; no margin-based claim \\
S5 & Selective replay and judge heterogeneity invalidate the archived
139/180 construction & mitigated for reporting: failure\_replay\_only;
top-up bounds reported; open experimentally: homogeneous rerun \\
S6 & Direction can invert on LongMemEval (B4 63.3\% \textgreater{}
controller 61.1\%) & open: no comparative claim until rerun \\
\end{longtable}
}

\subsection{9.5 MGP-6 v4.3-specific
threats}\label{mgp-6-v4.3-specific-threats}

The v4.3 release repairs the v2 construct gate, but it introduces its
own threats and none is closed by a higher accuracy number.

{\def\LTcaptype{none} % do not increment counter
\begin{longtable}[]{@{}
  >{\raggedright\arraybackslash}p{(\linewidth - 6\tabcolsep) * \real{0.1397}}
  >{\raggedright\arraybackslash}p{(\linewidth - 6\tabcolsep) * \real{0.2866}}
  >{\raggedright\arraybackslash}p{(\linewidth - 6\tabcolsep) * \real{0.2866}}
  >{\raggedright\arraybackslash}p{(\linewidth - 6\tabcolsep) * \real{0.2870}}@{}}
\toprule\noalign{}
\begin{minipage}[b]{\linewidth}\raggedright
\textbf{\#}
\end{minipage} & \begin{minipage}[b]{\linewidth}\centering
\textbf{Threat}
\end{minipage} & \begin{minipage}[b]{\linewidth}\centering
\textbf{Instance}
\end{minipage} & \begin{minipage}[b]{\linewidth}\centering
\textbf{Status}
\end{minipage} \\
\midrule\noalign{}
\endhead
\bottomrule\noalign{}
\endlastfoot
V1 & No hidden comparator. Only one configuration was authorised on the
sealed split & every baseline, ablation, and external-system contrast in
§7.10.6 is development evidence on the split that selected the
configuration & open: hidden baselines must be scored under the same
one-shot discipline before any comparative claim reaches the sealed
split \\
V2 & Development intervals are too narrow. Component-cluster resampling
is unavailable to the public scorer & §7.10.6 intervals are item
bootstraps; the paired McNemar tests assume independence that does not
hold within a component & partially mitigated: relied upon only for
contrasts with b/c splits of 98/6 or larger; the 0.117 successor gain is
flagged as the marginal case \\
V3 & Annotation provenance is attested, not authenticated & dual packets
agree on all 900 items (κ = 1.0), but the builder validates records and
hashes rather than reviewer identity or duration & open: perfect
agreement is also consistent with a manual that is unambiguous by
construction, so it certifies followability, not construct validity \\
V4 & Abstention is not calibrated selection & selective accuracy 0.8449
\textless{} overall accuracy 0.8481 & disclosed, unrepaired: the guard
performs governance, not selective-risk targeting; the two must not be
conflated \\
V5 & Successor configuration is untested on any sealed split & 0.9611
development accuracy, guard-composed & open: must not be read as an
improvement over the 0.8481 hidden result --- different splits,
different exposure \\
V6 & Cost is unmeasured in currency & provider telemetry emitted no
price field on any run & mitigated by disclosure: tokens and latency
reported as measured; no list-price reconstruction \\
V7 & Single model family, single seed & all model-backed rows use
qwen/qwen3.5-9b, seed 4301, temperature 0 & open: intervals reflect item
and component sampling, not run or model-family variance \\
V8 & Latency is provider wall time, not compute & measured through a
shared OpenRouter route with concurrency 1 & mitigated by framing:
reported as operational cost, never as a compute or throughput claim \\
\end{longtable}
}

\subsection{9.6 Priority actions (decreasing validity
impact)}\label{priority-actions-decreasing-validity-impact}

\begin{enumerate}
\def\labelenumi{\arabic{enumi}.}
\item
  Execute the role-separated MGP-6 v3 confirmation protocol: independent
  writers, all actions within every family, long trajectories, sealed
  labels, and shortcut rejection gates.
\item
  Recruit three blinded human annotators for the confirmation test and
  retain raw labels plus independent adjudication.
\item
  Evaluate component-oracle and predicted-component pipelines with
  family-clustered, cost-sensitive analysis.
\item
  Run external memory systems with native observable metrics; treat
  action mappings only as sensitivity analyses.
\item
  Execute a complete homogeneous downstream run only after its input,
  provider, judge, and provenance gates pass.
\end{enumerate}

\subsection{9.7 Ethics, privacy, and dual
use}\label{ethics-privacy-and-dual-use}

A system that supersedes and audits memories can be abused to rewrite
records or manufacture provenance; a system that retains them enables
surveillance, extraction, and membership inference (Carlini et al.,
2021; Shokri et al., 2017; Bender et al., 2021; Weidinger et al., 2022).
Mitigations in this project\textquotesingle s scope: lineage is
append-only and superseded records are preserved rather than destroyed;
ABSTAIN routes ambiguity to review instead of committing; the evaluation
data is synthetic (MGP-6) or externally licensed and not redistributed
(LongMemEval); no personal data of real users is stored. Deployment
beyond research requires access control, dual control for bulk
operations, and organizational governance consonant with risk-management
frameworks (Standards and Technology, 2023); alignment techniques at the
model layer (Ouyang et al., 2022; Bai et al., 2022) complement but do
not replace the deterministic write boundary. The governance controller
is an engineering framework, not legal advice; "supersede" and "abstain"
here are epistemic operations, not compliance guarantees.

\subsection{9.8 Reproducibility
statement}\label{reproducibility-statement}

Released: code, MGP-6 v2 data, per-item predictions, the 10,000-sample
construct-validity audit, cluster intervals, cost sensitivity,
bibliographic checks, and the v3 confirmation protocol. The v4.3 final
configuration, freeze authorization, aggregate hidden metrics, and
prediction hash are reproducibly recorded, while hidden records and gold
remain private. The few-shot run is incomplete; mapped external systems
and all LongMemEval evidence are non-primary. Independent v3
construction and three-human annotation remain unexecuted. Documentation
follows datasheet/model-card practice (Gebru et al., 2021; Mitchell et
al., 2019).

\section{10. Conclusion and Future
Work}\label{conclusion-and-future-work}

\subsection{10.1 Conclusion}\label{conclusion}

This thesis treats the memory write as the governed object: an explicit
KEEP/SUPERSEDE/ABSTAIN decision over typed, provenance-tracked,
lineage-preserving records, measured separately at the action,
stored-state/lineage, and answer layers. It reaches two conclusions, one
methodological and one empirical, and they arrived in that order.

The methodological conclusion is negative and it came first. MGP-6 v2
fails a construct-validity gate: conflict type alone reaches 1.000
held-out accuracy without reading the dialogue or the memory state. On
that pilot the lexical reference policy obtains 0.8889 test accuracy
against 0.9000 for the structured LLM with a paired family interval
containing zero; the reference has 0/33 observed action false commits
but a one-sided upper bound of 0.0868 and 10/45 missed updates, and the
cost matrix reverses which policy is preferable. No capability claim
survives that gate, and the thesis makes none.

The empirical conclusion follows from building the instrument the
failure demanded. MGP-6 v4.3 --- 900 human-authored items in 100
semantic components across 20 domains, component-split, exactly
action-balanced, dual-annotated, with the shortcut probe enforced as a
release gate --- supports a single pre-registered configuration scored
once on a sealed 540-item split: \textbf{0.8481 accuracy, 0.8479
macro-F1, Cohen\textquotesingle s κ 0.7722}, with component-cluster
intervals roughly 40\% wider than item independence would imply. Against
a \$1/3\$ chance floor, every structural baseline, every component
ablation, the zero-shot policy, the standalone rule policy, and three
deployed external memory systems fall between 0.2944 and 0.3667 on
development data, failing in two opposite directions --- commit-happy
systems overwrite 77--98\% of records they should keep, abstain-happy
systems decline 78--98\% of decisions they should make. Downstream, a
homogeneous full-slice LongMemEval-M run places governed writes at
76.1\% above every structural baseline, with the gain concentrated in
knowledge-update and preference competencies rather than spread
uniformly.

The thesis\textquotesingle s most transferable finding is smaller than
any of those numbers and more useful than all of them: \textbf{for a
write-governance system the error structure carries more information
than the aggregate}. Of 38 mishandled ambiguous items, 30 leave the
stored state byte-identical and lose only the audit record, while 8
replace an active record on insufficient evidence. A single false-commit
rate of 0.2111 overstates the state-corrupting harm fivefold. The same
discipline exposes a defect the aggregate would have hidden: selective
accuracy (0.8449) sits below overall accuracy (0.8481), so the
abstention mechanism performs governance without performing
selective-risk targeting, and the two are not the same capability.

What the evidence does not support bounds all of it. Only one hidden
configuration was authorised, so every comparative contrast is
development evidence on the split used for selection. The collections
are synthetic, English-only, single-slot, and single-step. Human
annotation is owner-attested rather than identity-authenticated, and
perfect agreement is equally consistent with a labelling manual that is
unambiguous by construction. One model family, one seed. Monetary cost
is unavailable rather than estimated. What the thesis offers is an
auditable protocol, a benchmark whose construct gates run at build time,
and a measured result inside explicitly stated bounds --- not controller
superiority, population safety, or deployment performance.

\subsection{10.2 Answers to the research
questions}\label{answers-to-the-research-questions}

\textbf{RQ1.} A governed write path makes actions and state transitions
explicit and auditable, and on MGP-6 v4.3 the composed controller
performs the three-way decision far above chance (0.8481 sealed-hidden,
κ 0.7722) where every baseline, ablation, and deployed external system
does not. Component decomposition attributes the result to the
composition rather than to either part alone. It does not establish that
this generalizes beyond a controlled synthetic diagnostic, and no hidden
comparator exists.

\textbf{RQ2.} MGP-6 v2 exposed its own label shortcut and is suitable
only for mechanism debugging. MGP-6 v4.3 moves that probe to a
build-time release gate, so the successor instrument is fit for
controlled comparison --- but its components, balanced evidence, and
constructed scenarios still do not reproduce the distribution of natural
user interactions.

\textbf{RQ3.} Decision metrics, commit safety, state and lineage
metrics, probe behaviour, and judged answers must be reported separately
--- and, this thesis adds, so must the \emph{classes} of decision error.
Separating state-preserving from state-corrupting write failures changes
the estimated harm fivefold, and separating governance from selective
risk revealed that the controller\textquotesingle s abstentions are not
risk-targeted. Two shipped forgetting metrics fail and are disclosed as
failing.

\textbf{RQ4.} Component-cluster resampling, shortcut probes,
finite-sample bounds, and cost sensitivity materially change what may be
claimed. On v2 they turn a policy comparison into a null. On v4.3 they
widen the reported interval by about a third relative to item
independence, and they show the ASK controller ordering is stable across
all three cost regimes where the v2 comparison flipped.

\textbf{RQ5.} A frozen ledger with claim-eligibility labels, mandatory
phrasing rules, co-reported raw and adjudicated numbers, enumerated
invalid runs, and a recomputation artifact that regenerates every
reported number from raw per-item predictions keep thesis, papers, and
code consistent --- including where the consistent statement is
"development evidence only" or "not final".

\subsection{10.3 Future work}\label{future-work}

\textbf{Immediate --- close the hidden-comparator gap.} The single most
valuable next experiment is not a better controller but a \emph{scored
comparator set}: run the four structural baselines, the four ablations,
and the successor configuration against the sealed v4.3 hidden split
under the same one-shot pre-registration discipline that produced the
0.8481 result. Until that exists, §7.10.6 remains development evidence
and the successor\textquotesingle s 0.9611 is an unverified development
number.

\textbf{Immediate --- independent construction.} Execute
benchmarks/mgp6/MGP6\_V3\_INDEPENDENT\_PROTOCOL.md with writers
separated from controller design, three blinded human annotators whose
identity the build system can authenticate, all three actions inside
every conflict family, long trajectories, and pre-unblinding shortcut
gates. The v4.3 build-time gate is a real improvement over the v2
post-hoc audit, but it is still run by the project that authored the
data.

\textbf{Near-term --- calibrate the abstention.} The selective-accuracy
inversion (0.8449 versus 0.8481) identifies a concrete target: the guard
currently decides \emph{whether a write is justified}, not \emph{whether
the controller is likely to be wrong}. A confidence signal fitted on an
independent set, evaluated by risk--coverage rather than accuracy, would
separate these two functions and is testable on development data without
consuming another hidden run.

\textbf{Near-term --- variance the current design cannot see.} All
reported intervals reflect item and component sampling under one model
family, one seed, and temperature 0. Multiple seeds, multiple model
families, and repeated runs of the same configuration would quantify run
variance, which is currently unmeasured and is the most likely source of
an over-read on any single number in this thesis.

\textbf{Near-term --- the designed but unexecuted SSP intervention.} The
provenance-constrained typed preference record and answer contract
(SINGLE\_SESSION\_PREFERENCE\_METHOD.md), with its credit-attribution
ablation between grounding-supply and answer-side repairs, targets the
competency where the V12 controller retains the most headroom (0.533).

\textbf{Longer-term.} Evaluate oracle versus predicted components; run
external memory systems under their native operations with observable
state, stale retrieval, latency, and cost, treating action mappings only
as sensitivity analyses; execute a complete homogeneous downstream
pipeline on the full LongMemEval collection rather than a 180-item
slice; and connect the operational policy to non-prioritized belief
revision (Alchourrón et al., 1985) and multi-principal governance (Ren
et al., 2026).

The broader conclusion stands independent of the pending cells: a
memory-capable agent should be judged not only by what it recalls, but
by what it declines to commit, what it supersedes with auditable
lineage, and whether its evaluation can tell those apart --- and, this
thesis adds, by whether its evaluation can tell a lost audit record
apart from a corrupted one.

\section{Appendix A. Data and Record
Schemas}\label{appendix-a.-data-and-record-schemas}

\subsection{A.1 MGP-6 v2 scenario row (schema:
benchmarks/mgp6/mgp6\_schema\_v2.py; spec:
mgp6\_v2\_spec.md)}\label{a.1-mgp-6-v2-scenario-row-schema-benchmarksmgp6mgp6_schema_v2.py-spec-mgp6_v2_spec.md}

Fields (abridged): scenario\_id (stable ID, e.g. MGP6V2-T3-TEST-016),
type (T1--T6), split (train/dev/test), context
(prior-memory-establishing turns), candidate (new event), gold\_action
(KEEP/SUPERSEDE/ABSTAIN), query + query\_mode (direct\_recall /
premise\_resistance / implicit\_application), gold\_answer.

\subsection{A.2 Released prediction row
(abridged)}\label{a.2-released-prediction-row-abridged}

\{\\
"scenario\_id": "MGP6V2-T1-TEST-016",\\
"predicted\_action": "ABSTAIN",\\
"predicted\_memory\_state": {[}\\
\{"id": "mgp6v2-t1-test-016-m1", "status": "active", "supersedes\_id":
null\}\\
{]},\\
"predicted\_answer": "The user\textquotesingle s current desk lighting
preference is warm light.",\\
"confidence": 0.4\\
\}

This is an abridgement of the restored released artifact, not a
synthetic scorecard trace. Gold labels are joined from the dataset
during evaluation, and system/run identity is encoded by the artifact
path (for example,
artifacts/mgp6\_v2/ref\_5systems/controller/predictions.jsonl).

\subsection{A.3 Memory record and lineage (store:
src/ask/memory/operations.py)}\label{a.3-memory-record-and-lineage-store-srcaskmemoryoperations.py}

Records carry: id, type (e.g. stable\_pref / transient\_pref /
episodic\_fact), content, holder, provenance (source turn/session),
status (active/superseded), superseded\_by lineage link, timestamps.
retrieve\_active excludes superseded records; get\_lineage walks the
chain.

\section{Appendix B. Prompt Contracts
(pointer)}\label{appendix-b.-prompt-contracts-pointer}

All prompt templates are versioned under prompts/
(extraction\_v0\ldots v\_final3.j2, relation\_stage1\_v0.j2,
relation\_stage2\_v0.j2) and referenced by the run configs. Contract
requirements: strict JSON output schemas; retrieved memory delimited as
untrusted data; the action vocabulary enumerated; unparsable output
resolves to no mutation. Judge prompts and guard order are documented in
src/ask/eval/longmemeval.py (§4.8). Prompts are not duplicated here to
keep a single source of truth.

\section{Appendix C. Metric
Definitions}\label{appendix-c.-metric-definitions}

Let yᵢ ∈ \{K, S, A\} be the gold action and ŷᵢ the canonicalized
prediction (unknown → INVALID, always incorrect).

\begin{itemize}
\item
  \textbf{Action accuracy} = 1/nΣᵢ 1{[}yᵢ = ŷᵢ{]}; \textbf{macro-F1} =
  unweighted mean of per-action F1.
\item
  \textbf{Invalid-commit rate} = Pr(ŷᵢ ∈ \{K, S\} \textbar{} yᵢ = A) ---
  commits on gold-ABSTAIN items (definition per revised
  metrics\_mgp6.py).
\item
  \textbf{Lineage integrity}: share of supersession chains whose links
  are complete and correct after the run.
\item
  \textbf{Stale-retrieval rate}: share of probes on which a superseded
  record surfaces as current.
\item
  \textbf{Premise resistance}: share of premise\_resistance probes on
  which the stale/false premise is rejected.
\item
  \textbf{Implicit application}: share of implicit\_application probes
  answered from the current value.
\item
  \textbf{LME answer accuracy}: judge-graded correctness; per-competency
  accuracy with stratified percentile-bootstrap CIs (SSP: n=20,000, seed
  0).
\item
  \textbf{Primary statistical machinery}: 10,000-replicate paired
  cluster bootstrap over template\_family\_id, repeated over
  minimal-pair groups plus singleton items; family-macro and
  item-weighted estimates; one-sided finite-sample upper bounds for zero
  event counts; frozen asymmetric expected-cost matrices; descriptive
  ECE/Brier and risk--coverage diagnostics. Exact item-level McNemar,
  Holm--Bonferroni, Newcombe intervals, and Wilson marginals are
  retained only as secondary historical analyses because item
  independence is doubtful.
\item
  \textbf{Excluded}: shipped over\_forgetting\_rate /
  under\_forgetting\_rate (saturate at 1.0 for every system; ledger
  known\_metric\_issues).
\end{itemize}

\section{Appendix D. Annotation Manual and Adjudication
(pointer)}\label{appendix-d.-annotation-manual-and-adjudication-pointer}

The annotation manual is benchmarks/mgp6/annotation\_manual\_v2.md; the
kit builder is make\_annotation\_kit.py; agreement computation is
compute\_agreement\_v2.py; raw annotation files live under
benchmarks/mgp6/annotations/human/ (independent annotator) and
mgp6\_v2/annotation/ (four-LLM panel). The single open adjudication item
is MGP6V2-T3-TEST-016 (gold SUPERSEDE vs annotator ABSTAIN, implicit
conflict). Adjudication protocol: blinded display of the raw item and
both labels; evidence spans and policy principle for each; resolution ∈
\{released, alternative, both defensible, underspecified\}; logged with
adjudicators, date, and version impact. Main results use released gold;
sensitivity may score the alternative label or exclude the item (±1 item
= 1.1\% of the test split).

\section{Appendix E. Reproducibility
Commands}\label{appendix-e.-reproducibility-commands}

\# Environment and offline checks\\
make test \# unit + regression tests\\
python benchmarks/mgp6/validate\_mgp6.py \# dataset integrity\\
\strut \\
\# MGP-6 systems and ablations (see benchmarks/mgp6/EVAL\_HARNESS.md)\\
python benchmarks/mgp6/run\_5systems.py\\
python benchmarks/mgp6/run\_ablations.py\\
python benchmarks/mgp6/run\_llm\_controller.py\\
python benchmarks/mgp6/run\_llm\_decision\_baseline.py\\
\strut \\
\# Statistics (stdlib-only, read-only on run artifacts)\\
python3 scripts/analysis/audit\_mgp6\_validity.py\\
\# writes artifacts/analysis/mgp6\_validity\_audit.\{json,md\}\\
python3 scripts/analysis/mcnemar\_mgp6.py \# prints table + JSON\\
\# item-level historical output:
reports/statistical\_analysis\_raw\_output.txt\\
\strut \\
\# Citation coverage and metadata quality\\
python3 scripts/release/audit\_bibliography.py\\
\strut \\
\# SSP analysis\\
python scripts/evaluation/evaluate\_single\_session\_preferences.py\\
\strut \\
\# Final tables (deterministic transcription of the ledger)\\
python scripts/analysis/build\_final\_results.py\\
\strut \\
\# LongMemEval final-1.1 protocol (EXTERNALLY PENDING)\\
\# configs/release/final.yaml + FINAL\_EVALUATION\_PROTOCOL.md\\
\# The supplied OpenRouter credential returned HTTP 403 before
inference.

Credentials are supplied only via environment variables; no credential
appears in configs, logs, or this repository.

\section{Appendix F. Artifact Manifest and Claim-Eligibility
Vocabulary}\label{appendix-f.-artifact-manifest-and-claim-eligibility-vocabulary}

{\def\LTcaptype{none} % do not increment counter
\begin{longtable}[]{@{}
  >{\raggedright\arraybackslash}p{(\linewidth - 4\tabcolsep) * \real{0.2774}}
  >{\raggedright\arraybackslash}p{(\linewidth - 4\tabcolsep) * \real{0.3612}}
  >{\raggedright\arraybackslash}p{(\linewidth - 4\tabcolsep) * \real{0.3614}}@{}}
\toprule\noalign{}
\begin{minipage}[b]{\linewidth}\raggedright
\textbf{Artifact}
\end{minipage} & \begin{minipage}[b]{\linewidth}\centering
\textbf{Path}
\end{minipage} & \begin{minipage}[b]{\linewidth}\centering
\textbf{Role}
\end{minipage} \\
\midrule\noalign{}
\endhead
\bottomrule\noalign{}
\endlastfoot
Frozen run ledger & artifacts/results\_ledger.json v1.1 (sha256
fea886f1ae746e921e785536e6ce47bea89ec11b99f8f2e75936dc71c6c5185e) &
historical run provenance and archived result transcription; its older
policy interpretation is superseded \\
Construct-validity audit &
artifacts/analysis/mgp6\_validity\_audit.\{json,md\} & primary source
for MGP-6 shortcut probes, held-out policy metrics, cluster intervals,
calibration, finite-sample bounds, and cost sensitivity \\
Final tables & reports/FINAL\_RESULTS.\{md,json\} & deterministic ledger
transcription \\
Statistics & reports/STATISTICAL\_ANALYSIS.md,
reports/statistical\_analysis\_raw\_output.txt & McNemar/Holm/CIs \\
Claims control & reports/CLAIMS\_EVIDENCE\_MATRIX.md; thesis-scoped:
thesis/claims\_evidence\_matrix.md & phrasing rules, strength tags \\
Source of truth & docs/audit/SCIENTIFIC\_SOURCE\_OF\_TRUTH.md & claim
strengths, corrections \\
Error analysis & reports/ERROR\_TAXONOMY.md;
docs/research/SINGLE\_SESSION\_PREFERENCE\_FAILURE\_ANALYSIS.md;
\ldots/SINGLE\_SESSION\_PREFERENCE\_METHOD.md & Chapter 8 sources \\
Validity & reports/THREATS\_TO\_VALIDITY.md;
docs/audit/DATA\_LEAKAGE\_AND\_RESEARCHER\_DOF.md & Chapter 9 sources \\
Run reconstruction & docs/audit/CONTROLLERRUN\_RECONSTRUCTION.md;
\ldots/CONTROLLERRUN\_FAILURE\_LEDGER.md;
\ldots/INVALID\_AND\_SUPERSEDED\_RUNS.md & LongMemEval anatomy \\
Related work audit & reports/RELATED\_WORK\_UPDATE.md & Chapter 2
source \\
Dataset + kit & benchmarks/mgp6/ (data, manual, metrics, annotation) &
Chapter 5 \\
MGP-6 run artifacts & artifacts/mgp6\_v2/\ldots{} (per-item predictions
+ reports) & Chapter 7 \\
LME-M run evidence & external run pointers in
docs/audit/RUN\_CLASSIFICATION.csv & Chapters 7--8; raw external data
not redistributed \\
Reproducibility & reports/REPRODUCIBILITY\_CHECKLIST.md, Makefile,
.github/workflows/ci.yml & Appendix E \\
\end{longtable}
}

\textbf{Claim-eligibility vocabulary:} diagnostic\_only --- informative
but not independent validation; failure\_replay\_only --- selected by
prior failure and excluded from result claims; incomplete --- missing
required coverage; invalid --- unusable; externally\_pending ---
specified but blocked by an external prerequisite. The older
final\_candidate label is a provenance state, not permission to ignore
the later shortcut failure.

\section{References}\label{references}

Citation keys resolve against thesis/references.bib (single consolidated
bibliography; every cited key was rechecked on 2026-07-15 and is
present). Metadata corrections use primary publication records where
available; unresolved institutional formatting is kept separate from
bibliographic identity. The final bibliography style will be selected
after the official University of Tehran template is supplied.

\subsection{Compiled reference list}\label{compiled-reference-list}

Alchourrón, Carlos E.; Gärdenfors, Peter; Makinson, David (1985). On the
Logic of Theory Change: Partial Meet Contraction and Revision Functions.
Journal of Symbolic Logic. DOI: 10.2307/2274239.

Bai, Yuntao; Kadavath, Saurav; Kundu, Sandipan; Askell, Amanda; Kernion,
Jackson; Jones, Andy; Chen, Anna; Goldie, Anna; Mirhoseini, Azalia;
McKinnon, Cameron; others (2022). Constitutional AI: Harmlessness from
AI Feedback. arXiv preprint arXiv:2212.08073.

Bender, Emily M.; Gebru, Timnit; McMillan-Major, Angelina; Shmitchell,
Shmargaret (2021). On the Dangers of Stochastic Parrots: Can Language
Models Be Too Big?. Proceedings of the 2021 ACM Conference on Fairness,
Accountability, and Transparency. DOI: 10.1145/3442188.3445922.

Borgeaud, Sebastian; Mensch, Arthur; Hoffmann, Jordan; Cai, Trevor;
Rutherford, Eliza; Millican, Katie; van den Driessche, George; Lespiau,
Jean-Baptiste; Damoc, Bogdan; Clark, Aidan; de Las Casas, Diego; Guy,
Aurelia; Menick, Jacob; Ring, Roman; Hennigan, Tom; Huang, Saffron;
Maggiore, Loren; Jones, Chris; Cassirer, Albin; Brock, Andy; Irving,
Geoffrey; Vinyals, Oriol; Osindero, Simon; Simonyan, Karen; Rae, Jack
W.; Elsen, Erich; Sifre, Laurent (2022). Improving Language Models by
Retrieving from Trillions of Tokens. Proceedings of the 39th
International Conference on Machine Learning.

Carlini, Nicholas; Tramèr, Florian; Wallace, Eric; Jagielski, Matthew;
Herbert-Voss, Ariel; Lee, Katherine; Roberts, Adam; Brown, Tom; Song,
Dawn; Erlingsson, Úlfar; Oprea, Alina; Raffel, Colin (2021). Extracting
Training Data from Large Language Models. 30th USENIX Security
Symposium.

Chao, Hanxiang; Bai, Yihan; Sheng, Rui; Li, Tianle; Sun, Yushi (2026).
STALE: Can LLM Agents Know When Their Memories Are No Longer Valid?.
arXiv preprint arXiv:2605.06527. DOI: 10.48550/arXiv.2605.06527.

Chen, Ding; Niu, Simin; Li, Kehang; Liu, Peng; Zheng, Xiangping; Tang,
Bo; Li, Xinchi; Xiong, Feiyu; Li, Zhiyu (2025). HaluMem: Evaluating
Hallucinations in Memory Systems of Agents. arXiv preprint
arXiv:2511.03506. DOI: 10.48550/arXiv.2511.03506.

Chhikara, Prateek; Khant, Dev; Aryan, Saket; Singh, Taranjeet; Yadav,
Deshraj (2025). Mem0: Building Production-Ready AI Agents with Scalable
Long-Term Memory. arXiv preprint arXiv:2504.19413.

Chow, C. K. (1970). On Optimum Recognition Error and Reject Tradeoff.
IEEE Transactions on Information Theory. DOI: 10.1109/TIT.1970.1054406.

Cohen, Jacob (1960). A Coefficient of Agreement for Nominal Scales.
Educational and Psychological Measurement. DOI:
10.1177/001316446002000104.

Efron, Bradley; Tibshirani, Robert J. (1993). An Introduction to the
Bootstrap. Chapman and Hall/CRC.

El-Yaniv, Ran; Wiener, Yair (2010). On the Foundations of Noise-free
Selective Classification. Journal of Machine Learning Research.

Gebru, Timnit; Morgenstern, Jamie; Vecchione, Briana; Vaughan, Jennifer
Wortman; Wallach, Hanna; Daumé III, Hal; Crawford, Kate (2021).
Datasheets for Datasets. Communications of the ACM. DOI:
10.1145/3458723.

Geifman, Yonatan; El-Yaniv, Ran (2017). Selective Classification for
Deep Neural Networks. Advances in Neural Information Processing Systems.

Ha, Hyeonjeong; Kim, Jeonghwan; Qian, Cheng; Liu, Jiayu; Campbell,
William M.; Wu, Yue; Zhang, Yuji; McKeown, Kathleen; Hakkani-Tür, Dilek;
Ji, Heng (2026). MemGuard: Preventing Memory Contamination in Long-Term
Memory-Augmented Large Language Models. arXiv preprint arXiv:2605.28009.
DOI: 10.48550/arXiv.2605.28009.

Iscan, Mehmet (2026). Learning When to Remember: Risk-Sensitive
Contextual Bandits for Abstention-Aware Memory Retrieval in LLM-Based
Coding Agents. arXiv preprint arXiv:2604.27283. DOI:
10.48550/arXiv.2604.27283.

Jiang, Bowen; Hao, Zhuoqun; Cho, Young-Min; Li, Bryan; Yuan, Yuan; Chen,
Sihao; Ungar, Lyle; Taylor, Camillo J.; Roth, Dan (2025). Know Me,
Respond to Me: Benchmarking LLMs for Dynamic User Profiling and
Personalized Responses at Scale. Conference on Language Modeling (COLM).

Lewis, Patrick; Perez, Ethan; Piktus, Aleaskndra; Petroni, Fabio;
Karpukhin, Vladimir; Goyal, Naman; Kuttler, Heinrich; Lewis, Mike; Yih,
Wen-tau; Rocktäschel, Tim; Riedel, Sebastian; Kiela, Douwe (2020).
Retrieval-Augmented Generation for Knowledge-Intensive NLP Tasks.
Advances in Neural Information Processing Systems.

Li, Zhiyu; Xi, Chenyang; Li, Chunyu; Chen, Ding; Chen, Boyu; Song,
Shichao; Niu, Simin; Wang, Hanyu; Yang, Jiawei; Tang, Chen; Yu,
Qingchen; Zhao, Jihao; Wang, Yezhaohui; Liu, Peng; Lin, Zehao; Wang,
Pengyuan; Huo, Jiahao; Chen, Tianyi; Chen, Kai; Li, Kehang; Tao, Zhen;
Lai, Huayi; Wu, Hao; Tang, Bo; Wang, Zhengren; Fan, Zhaoxin; Zhang,
Ningyu; Zhang, Linfeng; Yan, Junchi; Yang, Mingchuan; Xu, Tong; Xu, Wei;
Chen, Huajun; Wang, Haofen; Yang, Hongkang; Zhang, Wentao; Xu, Zhi-Qin
John; Chen, Siheng; Xiong, Feiyu (2025). MemOS: A Memory OS for AI
System. arXiv preprint arXiv:2507.03724. DOI: 10.48550/arXiv.2507.03724.

Liu, Nelson F.; Lin, Kevin; Hewitt, John; Paranjape, Ashwin; Bevilacqua,
Michele; Petroni, Fabio; Liang, Percy (2024). Lost in the Middle: How
Language Models Use Long Contexts. Transactions of the Association for
Computational Linguistics. DOI: 10.1162/tacl\_a\_00638.

Madras, David; Pitassi, Toniann; Zemel, Richard (2018). Predict
Responsibly: Improving Fairness and Accuracy by Learning to Defer.
Advances in Neural Information Processing Systems.

Maharana, Adyasha; Lee, Dong-Ho; Tulyakov, Sergey; Bansal, Mohit;
Barbieri, Francesco; Fang, Yuwei (2024). Evaluating Very Long-Term
Conversational Memory of LLM Agents. Proceedings of the 62nd Annual
Meeting of the Association for Computational Linguistics (Volume 1: Long
Papers). DOI: 10.18653/v1/2024.acl-long.747.

Meng, Kevin; Bau, David; Andonian, Alex; Belinkov, Yonatan (2022).
Locating and Editing Factual Associations in GPT. Advances in Neural
Information Processing Systems.

Meng, Kevin; Sharma, Arnab Sen; Andonian, Alex; Belinkov, Yonatan; Bau,
David (2023). Mass-Editing Memory in a Transformer. International
Conference on Learning Representations (ICLR).

Mitchell, Margaret; Wu, Simone; Zaldivar, Andrew; Barnes, Parker;
Vasserman, Lucy; Hutchinson, Ben; Spitzer, Elena; Raji, Inioluwa
Deborah; Gebru, Timnit (2019). Model Cards for Model Reporting.
Proceedings of the Conference on Fairness, Accountability, and
Transparency. DOI: 10.1145/3287560.3287596.

Mozannar, Hussein; Sontag, David (2020). Consistent Estimators for
Learning to Defer to an Expert. Proceedings of the 37th International
Conference on Machine Learning.

Ouyang, Long; Wu, Jeff; Jiang, Xu; Almeida, Diogo; Wainwright, Carroll
L.; Mishkin, Pamela; Zhang, Chong; Agarwal, Sandhini; Slama, Katarina;
Ray, Alex; others (2022). Training Language Models to Follow
Instructions with Human Feedback. Advances in Neural Information
Processing Systems.

Packer, Charles; Wooders, Sarah; Lin, Kevin; Fang, Vivian; Patil,
Shishir G.; Stoica, Ion; Gonzalez, Joseph E. (2023). MemGPT: Towards
LLMs as Operating Systems. arXiv preprint arXiv:2310.08560.

Park, Joon Sung; O\textquotesingle Brien, Joseph; Cai, Carrie Jun;
Morris, Meredith Ringel; Liang, Percy; Bernstein, Michael S. (2023).
Generative Agents: Interactive Simulacra of Human Behavior. Proceedings
of the 36th Annual ACM Symposium on User Interface Software and
Technology. DOI: 10.1145/3586183.3606763.

Patel, Vedant (2026). Supersede: Diagnosing and Training the
Memory-Update Gap in LLM Agents. arXiv preprint arXiv:2606.27472. DOI:
10.48550/arXiv.2606.27472.

Rasmussen, Preston; Paliychuk, Pavlo; Beauvais, Travis; Ryan, Jack;
Chalef, Daniel (2025). Zep: A Temporal Knowledge Graph Architecture for
Agent Memory. arXiv preprint arXiv:2501.13956.

Reddy, Vikas; Challaram, Sumanth (2026). Don\textquotesingle t Ask the
LLM to Track Freshness: A Deterministic Recipe for Memory Conflict
Resolution. arXiv preprint arXiv:2606.01435. DOI:
10.48550/arXiv.2606.01435.

Ren, Zhe; Yang, Yibo; Chen, Yimeng; Zhao, Zijun; Fu, Benshuo; Shu,
Zhihao; Zhang, Bingjie; Xu, Yangyang; Guo, Dandan; Yan, Shuicheng
(2026). GateMem: Benchmarking Memory Governance in Multi-Principal
Shared-Memory Agents. arXiv preprint arXiv:2606.18829. DOI:
10.48550/arXiv.2606.18829.

Salemi, Alireza; Mysore, Sheshera; Bendersky, Michael; Zamani, Hamed
(2023). LaMP: When Large Language Models Meet Personalization. arXiv
preprint arXiv:2304.11406.

Shokri, Reza; Stronati, Marco; Song, Congzheng; Shmatikov, Vitaly
(2017). Membership Inference Attacks Against Machine Learning Models.
2017 IEEE Symposium on Security and Privacy. DOI: 10.1109/SP.2017.41.

Sun, H.; Zhang, Z.; Zeng, S. (2025). Preference-Aware Memory Update for
Long-Term LLM Agents. arXiv preprint arXiv:2510.09720.

Tan, Zhen; Yan, Jun; Hsu, I-Hung; Han, Rujun; Wang, Zifeng; Le, Long;
Song, Yiwen; Chen, Yanfei; Palangi, Hamid; Lee, George; Iyer, Anand
Rajan; Chen, Tianlong; Liu, Huan; Lee, Chen-Yu; Pfister, Tomas (2025).
In Prospect and Retrospect: Reflective Memory Management for Long-term
Personalized Dialogue Agents. Proceedings of the 63rd Annual Meeting of
the Association for Computational Linguistics (Volume 1: Long Papers).
DOI: 10.18653/v1/2025.acl-long.413.

Wang, Ziming (2026). TOKI: A Bitemporal Operator Algebra for
Contradiction Resolution in LLM-Agent Persistent Memory. arXiv preprint
arXiv:2606.06240. DOI: 10.48550/arXiv.2606.06240.

Wei, Lei; Peng, Xiao; Dong, Xu; Xie, Niantao; Wang, Bin (2026). FadeMem:
Biologically-Inspired Forgetting for Efficient Agent Memory. arXiv
preprint arXiv:2601.18642. DOI: 10.48550/arXiv.2601.18642.

Weidinger, Laura; Mellor, John; Rauh, Maribeth; Griffin, Conor; Uesato,
Jonathan; Huang, Po-Sen; Cheng, Myra; Glaese, Mia; Balle, Borja;
Kasirzadeh, Atoosa; others (2022). Taxonomy of Risks Posed by Language
Models. Proceedings of the 2022 ACM Conference on Fairness,
Accountability, and Transparency. DOI: 10.1145/3531146.3533088.

Wu, Di; Wang, Hongwei; Yu, Wenhao; Zhang, Yuwei; Chang, Kai-Wei; Yu,
Dong (2025). LongMemEval: Benchmarking Chat Assistants on Long-Term
Interactive Memory. International Conference on Learning Representations
(ICLR).

Xu, Wujiang; Liang, Zujie; Mei, Kai; Gao, Hang; Tan, Juntao; Zhang,
Yongfeng (2025). A-MEM: Agentic Memory for LLM Agents. arXiv preprint
arXiv:2502.12110.

Yan, Sikuan; Yang, Xiufeng; Huang, Zuchao; Nie, Ercong; Ding, Zifeng;
Li, Zonggen; Ma, Xiaowen; Bi, Jinhe; Kersting, Kristian; Pan, Jeff Z.;
Schütze, Hinrich; Tresp, Volker; Ma, Yunpu (2025). Memory-R1: Enhancing
Large Language Model Agents to Manage and Utilize Memories via
Reinforcement Learning. arXiv preprint arXiv:2508.19828. DOI:
10.48550/arXiv.2508.19828.

Zhao, Siyan; Hong, Mingyi; Liu, Yang; Hazarika, Devamanyu; Lin, Kaixiang
(2025). Do LLMs Recognize Your Preferences? Evaluating Personalized
Preference Following in LLMs. International Conference on Learning
Representations (ICLR).

Zhong, Wanjun; Guo, Lianghong; Gao, Qiqi; Ye, He; Wang, Yanlin (2024).
MemoryBank: Enhancing Large Language Models with Long-Term Memory.
Proceedings of the AAAI Conference on Artificial Intelligence.

National Institute of Standards; Technology (2023). Artificial
Intelligence Risk Management Framework (AI RMF 1.0). DOI:
10.6028/NIST.AI.100-1.
